% !TeX program = pdflatex
% IEEE Access adaptation of the current main-SMPT manuscript.
% The class reuses TeX's \year name; PGF needs the original numeric primitive.
\let\TeXyear\year
\documentclass{ieeeaccess}
\let\AccessYear\year
\let\year\TeXyear
\usepackage[T1]{fontenc}
\usepackage[utf8]{inputenc}
\usepackage[english]{babel}
\usepackage{cite}
\usepackage{amsmath,amssymb,amsfonts}
\usepackage{graphicx}
% Permit wide figures at the bottom as well as the top of a text page.
\usepackage{stfloats}
\usepackage{tikz}
\usepackage{pgfplots}
\usepgfplotslibrary{groupplots,fillbetween}
\usetikzlibrary{patterns,positioning,fit,arrows.meta}
\usepackage{booktabs,array,float,xcolor}
% Review replacements: new text is royal blue; original text is retained.
\definecolor{RoyalBlue}{RGB}{65,105,225}
\newcommand{\michael}[2]{{\color{RoyalBlue}#1}}
% Yulian uses the same new/previous-text convention, with changes in red.
\newcommand{\yulian}[2]{{\color{red}#1}}
% Use the template blue as process CMYK; PGF also manages PDF colour resources.
\definecolor{accessblue}{cmyk}{1,0.3,0,0.2}
\definecolor{greycolor}{cmyk}{0,0,0,0.8}
\definecolor{grey}{cmyk}{0,0,0,0.1}
\usepackage{microtype}
\microtypesetup{expansion=false}
\IfFileExists{xurl.sty}{\usepackage{xurl}}{\usepackage{url}}
\urlstyle{same}
\DeclareUrlCommand\path{\urlstyle{same}}
\setlength{\emergencystretch}{2em}
\hyphenation{inter-community intra-community Gillespie}
\floatstyle{ruled}
\newfloat{algorithm}{tbp}{loa}
\floatname{algorithm}{Algorithm}
\newtheorem{proposition}{Proposition}
\newtheorem{corollary}{Corollary}
\newtheorem{remark}{Remark}
\newenvironment{proof}{\begin{IEEEproof}}{\end{IEEEproof}}
\newcommand{\ibar}{\bar{y}}
\newcommand{\sbar}{\bar{s}}
\newcommand{\Exp}{\mathrm{Exp}}
\newcommand{\Full}{\textup{Full}}
\newcommand{\TS}{\textup{TwoScale}}
\newcommand{\ind}{\mathbf{1}}
\newcommand{\eW}{e_{W}}
% Keep the supplied footer's appearance without its overfull zero-width boxes.
\makeatletter
\newcommand{\AccessOddFooter}{\hbox to\textwidth{{\footervolfont VOLUME\ \thevol, \theyear}\hfil{\footerpagefont\thepage}}}
\newcommand{\AccessEvenFooter}{\hbox to\textwidth{{\footerpagefont\thepage}\hfil{\footervolfont VOLUME\ \thevol, \theyear}}}
\let\AccessTitleStyle\ps@titlepage
\let\AccessHeadingStyle\ps@headings
\def\ps@titlepage{\AccessTitleStyle
  \def\@oddfoot{\raisebox{3pt}{\AccessOddFooter}}
  \def\@evenfoot{\raisebox{3pt}{\AccessEvenFooter}}}
\def\ps@headings{\AccessHeadingStyle
  \def\@oddfoot{\raisebox{8pt}{\AccessOddFooter}}
  \def\@evenfoot{\raisebox{8pt}{\AccessEvenFooter}}}
\makeatother
\pagestyle{headings}
% Keep float spacing compact while retaining the IEEE page dimensions.
\setlength{\dbltextfloatsep}{12pt plus 2pt minus 2pt}
\setlength{\dblfloatsep}{10pt plus 2pt minus 2pt}
\renewcommand{\dbltopfraction}{0.9}
\renewcommand{\dblfloatpagefraction}{0.75}
% Allow for the bullet and gap in the template's abstract/index-term boxes.
% The class aliases these two lengths; separate them before adjusting the box.
\newlength{\AccessAbstractWidth}
\setlength{\AccessAbstractWidth}{\titlewidth}
\addtolength{\AccessAbstractWidth}{-10pt}
\let\abstractwidth\AccessAbstractWidth
% BEGIN EDITABLE PLOT STYLE
% Shared plotting style; all text inherits the article font.
\providecommand{\ArrivalCellFont}{\fontsize{6}{7}\selectfont}
\pgfplotsset{compat=1.18,paper axis/.style={font=\scriptsize,label style={font=\scriptsize},title style={font=\small},tick label style={font=\scriptsize},legend style={font=\scriptsize,draw=none,fill=none},axis line style={black!65},tick style={black!65},grid=major,grid style={black!10},scaled ticks=false,unbounded coords=jump,clip=true}}
\definecolor{FullBlue}{HTML}{2878B5}
\definecolor{ScaleOrange}{HTML}{D45A44}
\definecolor{Size50}{HTML}{D73027}
\definecolor{Size100}{HTML}{FC8D59}
\definecolor{Size250}{HTML}{4575B4}
\definecolor{Size1000}{HTML}{313695}
\pgfplotsset{colormap={paperdiv}{rgb=(0.01960784,0.18823529,0.38039216);
rgb=(0.12725875,0.39584775,0.66874279);
rgb=(0.26274510,0.57647059,0.76470588);
rgb=(0.56647443,0.76870434,0.86851211);
rgb=(0.81960784,0.89803922,0.94117647);
rgb=(0.96908881,0.96647443,0.96493656);
rgb=(0.99215686,0.85882353,0.78039216);
rgb=(0.95455594,0.64175317,0.50572857);
rgb=(0.83921569,0.37647059,0.30196078);
rgb=(0.69227220,0.09227220,0.16770473);
rgb=(0.40392157,0.00000000,0.12156863);}}
\pgfplotsset{colormap={paperseq}{rgb=(1.00000000,1.00000000,0.89803922);
rgb=(1.00000000,0.97539408,0.77193387);
rgb=(0.99764706,0.92156863,0.63607843);
rgb=(0.99607843,0.84347559,0.46915802);
rgb=(0.99607843,0.73490196,0.28000000);
rgb=(0.99497116,0.59747789,0.15949250);
rgb=(0.93960784,0.47137255,0.09490196);
rgb=(0.84822760,0.35229527,0.03497116);
rgb=(0.72000000,0.26039216,0.01098039);
rgb=(0.55686275,0.19123414,0.01737793);
rgb=(0.40000000,0.14509804,0.02352941);}}
% Stream all ECDF vertices to PGF without retaining PGFPlots point metadata.
% Arguments: display-coordinate file, stroke colour, dash setting.
\newcommand{\DrawExactECDF}[3]{%
  \begin{pgfonlayer}{main}%
  \begin{pgfscope}%
    \pgftransformshift{\pgfplotsqpointdescriptionxy{0}{0}}%
    \pgfsetxvec{\pgfpoint{\pgfkeysvalueof{/pgfplots/width}}{0pt}}%
    \pgfsetyvec{\pgfpoint{0pt}{\pgfkeysvalueof{/pgfplots/height}}}%
    \pgfpathrectangle{\pgfpointorigin}{\pgfpoint{\pgfkeysvalueof{/pgfplots/width}}{\pgfkeysvalueof{/pgfplots/height}}}%
    \pgfusepath{clip}%
    \pgfsetstrokecolor{#2}%
    \pgfsetlinewidth{.8pt}%
    #3%
    \pgfplothandlerconstantlineto%
    \pgfplotxyfile{#1}%
    \pgfusepath{stroke}%
  \end{pgfscope}%
  \end{pgfonlayer}%
}
% END EDITABLE PLOT STYLE

\begin{document}
% Publication metadata are assigned by the publisher.
\history{Manuscript prepared for review.}
\doi{To be assigned}
\vol{XX}
\AccessYear{XXXX}
\title{Two-scale stochastic simulation of epidemics on community networks with aggregate inter-community coupling}
\author{\uppercase{Yulian Kuryliak}\authorrefmark{1} and \uppercase{Michael T. M. Emmerich}\authorrefmark{2}}
\address[1]{Institute of Computer Science and Information Technologies, Lviv Polytechnic National University, 12 Stepan Bandera Street, Lviv 79000, Ukraine (e-mail: yulian.a.kuryliak@lpnu.ua)}
\address[2]{Faculty of Information Technology, University of Jyv\"askyl\"a, P.O. Box 35 (Agora), FI-40014 Jyv\"askyl\"a, Finland (e-mail: michael.t.m.emmerich@jyu.fi)}
\markboth{Kuryliak and Emmerich: Two-scale stochastic simulation with aggregate inter-community coupling}{Kuryliak and Emmerich: Two-scale stochastic simulation with aggregate inter-community coupling}
\corresp{Corresponding author: Yulian Kuryliak (e-mail: yulian.a.kuryliak@lpnu.ua).}
\begin{abstract}
\michael{Detailed within-community contact networks may be available or reconstructed, while connections between communities are known only through aggregate data.}{Stochastic epidemic simulation often lacks information about individual contacts. Aggregate data on movements between communities may nevertheless be available. Internal contact networks can be approximately reconstructed from structural characteristics, while specific inter-community contacts remain unknown.} We study a two-scale stochastic simulation method (\TS{}) that combines detailed local Gillespie simulators with aggregate inter-community transmission rates, without explicitly representing inter-community contacts. We prove that the ideal two-scale process is equivalent to exact simulation on a surrogate network with uniform inter-community connections. We identify the mechanisms of discrepancy from the full-network model and establish that the first infection export is stochastically no later than in the full-network model under uniformly random bridge-node placement in the theoretical setting considered. We assess the method's accuracy by comparing it with full-network Gillespie simulation \michael{on synthetic susceptible--infected networks}{on synthetic SI networks}, including 48 configurations of chains of ten communities. The experiments showed close time distributions when infection spread rapidly within communities: in a chain of ten cliques, the normalised Wasserstein distance after nine inter-community crossings was $3.34\%$. However, discrepancies increased substantially in communities with longer internal paths and lower connectivity, reaching $40.9\%$ in the most discrepant configuration. Their magnitude also depended on community size and the number of inter-community edges. \michael{The ratio of within-community infection-time dispersion to single-edge transmission time showed a strong rank association with the discrepancies}{The proposed time-scale measure showed a strong rank association with the magnitude of the discrepancies} (Spearman correlation $0.91$). These results help identify conditions under which aggregate connections reproduce temporal characteristics close to full-network simulation without specifying individual inter-community contacts.
\end{abstract}
\begin{keywords}
Aggregate inter-community connections, Community networks, Epidemic modelling, Gillespie algorithm, Incomplete contact data, Multiscale simulation, Simulation accuracy, Stochastic simulation.
\end{keywords}
\titlepgskip=-21pt
\maketitle

\section{Introduction}
\label{sec:introduction}

Epidemics spread at more than one scale. Inside a community, infection may pass through a detailed network of contacts between people \cite{keeling2005,pastorsatorras2015,kiss2017}. Between communities, the available data are often much less detailed. For example, they may give only passenger numbers, commuting flows, or a weighted link between two places \cite{balcan2009,barbosa2018human}. We therefore need a model that can use a detailed network inside each community and simpler data between communities.

This situation is common in practice. A local contact network may be observed, reconstructed, or generated to match known properties such as clustering and mean shortest-path length \cite{eubank2004,danon2011}. In contrast, links between communities may be known only through transport, migration, commuting, or mobile-phone data \cite{barbosa2018human,chang2021}. Such data can tell us the total strength of connection between two communities, but not which people form the connection. We may therefore know the total coupling without knowing the bridge nodes or their infection states.

If every contact is known, a continuous-time Markov chain (CTMC) can describe the epidemic, and the Gillespie method can simulate it exactly \cite{gillespie1977exact,gibson2000}. Efficient methods are also available for temporal, large, and sparse networks \cite{vestergaard2015,kuryliak2025simulating}. However, these methods still need the individual contacts that may transmit infection. Metapopulation models use aggregate flows between communities, but usually also replace the detailed local network by counts of people in different states \cite{colizza2007reaction,colizza2008,balcan2010modeling,ajelli2010}. Our problem lies between these two cases: we want to keep node-level stochastic spread inside every community while using only aggregate information between communities.

The idea of a two-scale simulator was presented in an earlier conference paper \cite{kuryliak2025towards}. It runs a separate local simulator for each community. The macro layer uses the current community states and connection weights to generate infections arriving from another community. In this article, we describe the method in detail and study its approximation error. In particular, we examine what is lost when real cross-community edges and their endpoints are replaced by the total connection weight, and an imported infection is placed on a randomly selected susceptible node.

Our main question is: when is this aggregate connection sufficient to reproduce the quantity of interest in the full-network CTMC? We study infection arrival times, their distributions, and whether transmission is possible through the real contacts. Similar arrival times do not guarantee a possible transmission route: the two-scale model can infect another community while all real bridge nodes in the source community are still susceptible. To test this, we use synthetic networks in which every contact is known. \michael{Full-network Gillespie simulation receives the whole network, whereas the two-scale method receives only the local networks and the total weights of the cross-community connections. Synthetic networks are a deliberate experimental choice: they provide a known microscopic reference, permit controlled variation of structure, and enable analytical checks that make the source of discrepancies interpretable. Explicit generators, parameters, and random seeds also support reproducible comparisons. This study therefore evaluates the approximation under controlled conditions; it does not establish empirical validity for a particular epidemic.}{Full-network Gillespie simulation receives the whole network, whereas the two-scale method receives only the local networks and the total weights of the cross-community connections.}

This article makes four contributions.
\begin{enumerate}
\item We describe how the simulator works: each community is simulated separately, and a shared scheduler determines infections between them. We explain how it coordinates the local simulators on time intervals of fixed length (Section~\ref{sec:method}, Algorithm~\ref{alg:scheduler}).
\item We show which network the method simulates exactly. The ideal process is exact for a \emph{surrogate network with uniform inter-community connections}: for each connected community pair, every node in one community is connected to every node in the other, and the total coupling weight is distributed equally among these node pairs (Proposition~\ref{prop:surrogate}). We then identify the sources of error and analyse a single export of infection (Proposition~\ref{prop:ordering}).
\item We compare the two-scale method with full-network Gillespie simulation for transmission to one neighbouring community and for spread along chains of communities. We measure differences using the relative shift in the median time of first infection and the Wasserstein distance between the distributions of these times (Sections~\ref{sec:design} and~\ref{sec:results}).
\item We propose an empirical indicator of accuracy for our benchmarks. It compares the spread of infection times inside a community with the typical transmission time through one cross-community edge (Section~\ref{sec:res_descriptor}).
\end{enumerate}

We consider susceptible--infected (SI) dynamics on static synthetic networks and use one fixed synchronisation step. Section~\ref{sec:related} explains the relation to earlier methods. Section~\ref{sec:method} describes the simulator. Section~\ref{sec:analysis} explains what the aggregate interface preserves and loses. Sections~\ref{sec:design} and~\ref{sec:results} present the experiments and results. Section~\ref{sec:discussion} discusses use, cost, and limitations.

%======================================================================
\section{Related work}
\label{sec:related}

\paragraph{Simulation of stochastic events.}
Gillespie's direct method draws the next event and its time from the current event rates \cite{gillespie1977exact}. The next reaction method keeps a separate clock for each possible event \cite{gibson2000}. Both methods are exact for the CTMC they simulate. When an event rate changes with time, its rate can be accumulated until it reaches a random threshold \cite{ethier1986}. Anderson~\cite{anderson2007} used this idea for systems with time-dependent rates and delays. Our macro clock uses the same idea, but its rate changes because the local epidemic states change. Network simulators can save work by updating only the rates affected by the latest event \cite{vestergaard2015,kuryliak2025simulating}. Our local simulators use the implementation of Kuryliak et al.~\cite{kuryliak2025simulating}.

Some multiscale methods speed up chemical-reaction models by approximating reactions that occur much faster than others \cite{haseltine2002,cao2005slow}. Our method removes a different kind of detail. It keeps the local infection events, but does not know which individual cross-community edges are currently able to transmit. Even so, the relative speeds of local spread and cross-community transmission help explain its error.

\paragraph{Models with several communities.}
Metapopulation models represent each subpopulation by counts of people in different states and connect the subpopulations by mobility flows. They have been used to study global spread and the conditions for invasion \cite{colizza2007reaction,colizza2008,balcan2009,balcan2010modeling}. Detailed mobility data can also be combined with local transmission models \cite{chang2021}. Agent-based models instead represent individuals throughout the model \cite{eubank2004,kerr2021}. Ajelli et al.~\cite{ajelli2010} compared these two broad approaches. \michael{The method studied here keeps individual nodes and contacts inside each community, but replaces the unknown individual contacts between communities by one aggregate stochastic interface. Related hybrid epidemiological models combine agent-based and population-based descriptions across regions or scales \cite{bicker2025}. Reconstruction studies also show that group structure and temporal contact statistics affect epidemic estimates from incomplete contact data \cite{genois2015sampling}; our focus is the error caused by aggregating only the inter-community interface.}{The method studied here keeps individual nodes and contacts inside each community, but replaces the unknown individual contacts between communities by one aggregate stochastic interface.}

\paragraph{Two levels of mixing and co-simulation.}
Some epidemic models combine contacts in small local groups with uniform mixing at a wider level \cite{ball1997}. Our surrogate network with uniform inter-community connections is related to this idea, but uniform mixing is used only between connected pairs of communities, and its strength is set by their aggregate weight.

The implementation also has features of co-simulation. In co-simulation, separate subsystem solvers move forward on a shared time grid and exchange a limited amount of information \cite{gomes2018cosim}. Here, the local community simulators are those subsystems. They send their infected and susceptible counts to the macro layer. The macro layer uses these data to generate the time of a cross-community infection and select the receiving community and target node. The local simulators can therefore remain black boxes for the scheduler. In this implementation, the exchanged data are used to create random infection events.

%======================================================================
\section{The two-scale simulation method}
\label{sec:method}

The experiments use one implementation of the two-scale simulator, in which the local simulators exchange information with the macro scheduler at fixed time intervals. To understand its error, we also consider a process with the same rates but exact timing of every event. We call this theoretical description the \emph{ideal process}. It separates the error caused by changing the network from the additional error caused by the time grid. The ideal process is exact for the surrogate network with uniform inter-community connections defined later; this does not make it exact for the original physical network. Below, we explain the reference model, the available data, and the steps of the implemented simulator.

\subsection{Full-network reference model}
\label{sec:reference}

Let $G=(V,E)$ be the contact network. Its nodes are divided into $M$ non-overlapping communities:
\begin{equation}
\label{eq:partition}
V=V_1\cup V_2\cup\cdots\cup V_M,
\qquad N_i=|V_i|.
\end{equation}
Here, $V_i$ is the node set of community $i$, and $N_i$ is its size. We write $E_{ii}$ for contacts inside community $i$ and $E_{ij}$ for contacts from nodes in community $i$ to nodes in community $j$. We treat each contact as an oriented transmission channel. An undirected contact therefore gives two channels, one in each direction.

Every node is susceptible ($S$), infected ($I$), or removed ($R$). The full network state at time $t$ is $X(t)$, and the state of community $i$ is $X_i(t)$. Let $S_i(t)$ and $I_i(t)$ be the sets of susceptible and infected nodes in that community. Their sizes are $s_i(t)$ and $y_i(t)$, and their fractions are $\sbar_i(t)=s_i(t)/N_i$ and $\ibar_i(t)=y_i(t)/N_i$. Below, $x$ denotes a particular network state; $x^{v:I}$ is the same state after node $v$ becomes infected, and $x^{u:R}$ is the state after node $u$ is removed.

The full model follows two simple rules. A susceptible node $v$ receives infection from each infected neighbour $u$ at rate $\beta_{uv}$. An infected node is removed at rate $\gamma$. We can write both rules as
\begin{equation}
\label{eq:full_generator}
q_v(x)=\sum_{\substack{u:(u,v)\in E\\x_u=I}}\beta_{uv}
\quad (x_v=S),
\qquad r_u(x)=\gamma\quad (x_u=I).
\end{equation}
Here, $q_v$ is the total infection rate of node $v$, and $r_u$ is the removal rate of node $u$. For a non-susceptible node, $q_v=0$; for a non-infected node, $r_u=0$. A rate describes how often an event occurs: for a constant rate $q$, the mean waiting time is $1/q$. For a cross-community channel, we write $\beta_{uv}=\beta_{\mathrm{out}}w_{uv}$, where $w_{uv}$ is its contact weight. We call exact simulation of these full-network rates \Full{}. All experiments in this article use SI dynamics, so $\gamma=0$.

The waiting time between state transitions is exponential and is determined by the total rate of all events in the current state $x$, which is
\[
a(x)=\sum_v q_v(x)+\sum_u r_u(x).
\]
If $a(x)>0$, the waiting time to the next event is exponentially distributed with rate $a(x)$ and mean $1/a(x)$. Infection of node $v$ is selected with probability $q_v(x)/a(x)$, and removal of node $u$ with probability $r_u(x)/a(x)$. After the event, the rates are updated. If $a(x)=0$, no further event can occur.

\subsection{Aggregate interface}
\label{sec:interface}

The two-scale simulator, denoted \TS{}, knows every local graph $G_i=(V_i,E_{ii})$ and its rates. It does not know the individual cross-community channels in $E_{ij}$. For each ordered pair of communities, it knows only the total weight
\begin{equation}
\label{eq:weight}
W_{ij}=\sum_{(u,v)\in E_{ij}}w_{uv},
\qquad i\neq j.
\end{equation}
If all weights equal one, $W_{ij}$ is the number of directed channels from $V_i$ to $V_j$. In an application, this total may instead come from transport, migration, mobility, or contact data. The macro layer also receives the current susceptible and infected counts from each local simulator. It does not know which nodes are endpoints of real cross-community edges. Table~\ref{tab:information} gives the main difference, and Figure~\ref{fig:schematic} shows it visually.

\begin{figure*}[!tbp]
\color{black}
\centering
\begin{tikzpicture}[
  x=1cm,y=1cm,
  font=\scriptsize,
  >=Latex,
  person/.style={circle,draw=black!60,fill=blue!18,minimum size=4.2mm,inner sep=0pt},
  infected/.style={person,fill=orange!70},
  bridge/.style={person,very thick,draw=black},
  bridgeinf/.style={infected,very thick,draw=black},
  commbox/.style={draw=black!45,dashed,rounded corners=2pt,inner sep=3pt},
  macrobox/.style={draw=black!70,rounded corners=2pt,fill=black!5,align=center,inner sep=3pt},
  link/.style={draw=black!55},
  weak/.style={draw=black!24,line width=.25pt},
  flow/.style={-{Latex[length=1.7mm]},draw=black!65}
]
% (a) Full microscopic interface
\node[font=\bfseries] at (0,2.25) {(a) Full network};
\node[infected] (a1) at (-1.65,1.20) {};
\node[person]   (a2) at (-1.55,0.35) {};
\node[bridgeinf](a3) at (-0.75,0.82) {};
\node[person]   (a4) at (-0.75,1.55) {};
\node[bridge]   (b1) at (0.65,1.15) {};
\node[person]   (b2) at (1.45,1.55) {};
\node[person]   (b3) at (1.55,0.55) {};
\node[bridge]   (b4) at (0.70,0.30) {};
\draw[link] (a1)--(a2)--(a3)--(a1) (a1)--(a4)--(a3);
\draw[link] (b1)--(b2)--(b3)--(b4)--(b1) (b1)--(b3);
\draw[very thick] (a3)--(b1);
\draw[very thick] (a3)--(b4);
\node[commbox,fit=(a1)(a2)(a3)(a4)] {};
\node[commbox,fit=(b1)(b2)(b3)(b4)] {};
\node at (-1.2,-0.10) {$V_i$};
\node at (1.15,-0.10) {$V_j$};

% (b) Aggregate community network with local contacts retained
\node[font=\bfseries] at (5.0,2.25) {(b) Aggregate network};
\node[infected] (e1) at (3.35,1.20) {};
\node[person]   (e2) at (3.45,0.35) {};
\node[infected] (e3) at (4.25,0.82) {};
\node[person]   (e4) at (4.25,1.55) {};
\node[person]   (f1) at (5.65,1.15) {};
\node[person]   (f2) at (6.45,1.55) {};
\node[person]   (f3) at (6.55,0.55) {};
\node[person]   (f4) at (5.70,0.30) {};
\draw[link] (e1)--(e2)--(e3)--(e1) (e1)--(e4)--(e3);
\draw[link] (f1)--(f2)--(f3)--(f4)--(f1) (f1)--(f3);
\node[commbox,fit=(e1)(e2)(e3)(e4)] (agg_i) {};
\node[commbox,fit=(f1)(f2)(f3)(f4)] (agg_j) {};
\draw[-{Latex[length=2.5mm,width=2.5mm]},draw=black!65,line width=3pt]
      (agg_i.east) -- node[above=5pt,text=black] {$W_{ij}$} (agg_j.west);
\node at (3.8,-0.10) {$V_i$};
\node at (6.15,-0.10) {$V_j$};

% (c) Surrogate interface
\node[font=\bfseries] at (10.25,2.25) {(c) Surrogate network};
\node[infected] (c1) at (8.85,1.45) {};
\node[person]   (c2) at (8.75,0.55) {};
\node[person]   (c3) at (9.55,1.05) {};
\node[person]   (c4) at (9.45,0.20) {};
\node[person]   (d1) at (11.00,1.45) {};
\node[person]   (d2) at (10.90,0.55) {};
\node[person]   (d3) at (11.70,1.05) {};
\node[person]   (d4) at (11.60,0.20) {};
\foreach \u in {c1,c2,c3,c4}{\foreach \v in {d1,d2,d3,d4}{\draw[weak] (\u)--(\v);}}
\draw[link] (c1)--(c2)--(c4)--(c3)--(c1);
\draw[link] (d1)--(d2)--(d4)--(d3)--(d1);
\node[commbox,fit=(c1)(c2)(c3)(c4)] {};
\node[commbox,fit=(d1)(d2)(d3)(d4)] {};
\node[align=center] at (10.25,-0.35) {$W_{ij}/(N_iN_j)$ per\\cross-community channel};
\end{tikzpicture}
\caption{\textcolor{black}{Three views of the interface. (a) In the full network, infection can cross only through the shown physical edges and their bridge nodes. (b) Internal nodes and contacts are retained in each community. Individual inter-community contacts are represented by one wide directed edge between $V_i$ and $V_j$, labelled with their total weight $W_{ij}$. (c) The ideal two-scale process is equivalent to a process on a surrogate network with uniform inter-community connections: for each connected community pair, all ordered node pairs have the same contact weight. The connections together have total weight $W_{ij}$. Orange nodes are infected and blue nodes are susceptible.}}
\label{fig:schematic}
\end{figure*}

\begin{table*}[!t]
\color{black}
\centering
\caption{\textcolor{black}{Data and rules of the two models. In all experiments, \TS{} selects the target node uniformly among susceptible nodes.}}
\label{tab:information}
\small
\begin{tabular}{>{\raggedright\arraybackslash}p{0.42\linewidth}>{\raggedright\arraybackslash}p{0.23\linewidth}>{\raggedright\arraybackslash}p{0.25\linewidth}}
\toprule
Information or action & \Full{} & \TS{}\\
\midrule
Local contacts, rates, and node states & known & known by local simulators\\
Individual cross-community edges and bridge nodes & known & unknown\\
Total cross-community weight $W_{ij}$ & obtained from the edges & given directly\\
Cross-community event rate & uses states of real edge endpoints & uses $\ibar_i$, $\sbar_j$, and $W_{ij}$\\
Target of a cross-community infection & endpoint of the active edge & a susceptible node drawn by the target rule\\
Number of possible SI infections through one bundle & at most the number of its channels & at most $s_j$ target infections\\
\bottomrule
\end{tabular}
\end{table*}

\subsection{Local simulators}
\label{sec:local}

Each community has a separate Gillespie simulator for its local graph $G_i$. It applies the infection and removal rates in~\eqref{eq:full_generator} using only the local edges $E_{ii}$. The implementation updates only rates affected by the latest event, following Kuryliak et al.~\cite{kuryliak2025simulating}. Each local simulator is independent and has its own random-number stream.

The macro layer can advance a local simulator to a specified time. The simulator applies the local events up to that time. The macro layer can then import an infection by changing one selected susceptible node to $I$. This changes the local rates. The implementation keeps the unused part of the random threshold for the next local event and recalculates its waiting time using the new total rate.

To recalculate the time of the next local event after an import, the simulator retains the remaining random threshold. Let $E_{\mathrm{loc}}\sim\mathrm{Exp}(1)$ be the threshold drawn after the previous local event, and let $A_{\mathrm{loc}}$ be the integral of the total local rate from that event to the import time. After the import, the total local rate is $a_{\mathrm{new}}$. If $a_{\mathrm{new}}>0$, the remaining waiting time is
\[
\Delta t=
\frac{E_{\mathrm{loc}}-A_{\mathrm{loc}}}{a_{\mathrm{new}}}.
\]
Thus, an import changes the waiting time, but not the random threshold itself. If another import occurs first, the accumulated rate is updated and the waiting time is recalculated again. If $a_{\mathrm{new}}=0$, no local event can occur until the rate becomes positive.

The event type and its node are selected using the updated rates. A new local threshold is drawn only if there is no existing one. Each local clock has its own threshold, separate from the macro threshold described in Section~\ref{sec:timing}. For prescribed import times and target nodes, the local simulator follows the CTMC rules of its own network. The approximation comes from the way imports are generated and their target nodes are selected.

\subsection{Macro layer}
\label{sec:macro}

The macro layer represents each community by its current fractions and represents each directed connection by $W_{ij}$. Its rate for an infection from community $i$ to community $j$ is
\begin{equation}
\label{eq:macro_rate}
\lambda_{ij}(t)=\beta_{\mathrm{out}}W_{ij}\ibar_i(t)\sbar_j(t),
\qquad
\Theta(t)=\sum_{i\neq j}\lambda_{ij}(t).
\end{equation}
The factors have a direct meaning: $\beta_{\mathrm{out}}W_{ij}$ gives the total connection strength, $\ibar_i$ is the infected fraction in the source, and $\sbar_j$ is the susceptible fraction in the receiving community. The sum $\Theta$ is the rate of any cross-community infection. Calculating the waiting time to the next event when $\Theta$ varies with time requires special treatment (Section~\ref{sec:timing}).

When such an event occurs, first select the source and receiving communities $(i,j)$ with probability $\lambda_{ij}/\Theta$. Then choose one susceptible node uniformly in the receiving community $j$ and infect it. If there are 20 susceptible nodes, each has probability $1/20$. This selection does not use real bridge nodes, because the aggregate interface does not know them. In later formulas, we write this target probability as $\Pi_{ij}(v\mid x)=1/s_j$ for $v\in S_j$, where $x$ is the current network state. Thus, $\Pi$ simply denotes the probability of choosing a particular susceptible node. Other target rules can be used if bridge-node information is available.

This aggregate rate is generally an approximation to transmission through the real physical edges. For fixed community counts, it is the expected rate under uniformly random selection of the edge endpoints, before their particular identities, states, and epidemic history are taken into account. This does not guarantee the correct rate for an epidemic already realised on specific edges. For uniform all-to-all coupling, it is the actual rate, as explained in Section~\ref{sec:analysis}.

\subsection{Timing and synchronisation}
\label{sec:timing}

After a macro event at time $t_0$, draw a random threshold $\michael{\mathcal E}{E}\sim\Exp(1)$: an exponential random number with mean one. Starting from zero, accumulate the total macro rate $\Theta$:
\begin{equation}
\label{eq:hazard_integral}
A(t)=\int_{t_0}^{t}\Theta(s)\,ds,
\qquad A(t_0)=0.
\end{equation}
The next macro event occurs at the first time $t^\ast$ when $A(t^\ast)=\michael{\mathcal E}{E}$.

In words, the area under the rate curve grows until it first reaches the random threshold. If it never reaches the threshold, no further macro event occurs. This is the standard cumulative-rate construction \cite{ethier1986,anderson2007}. If $\Theta$ is constant, the waiting time is $\michael{\mathcal E}{E}/\Theta$. Drawing $\michael{\mathcal E}{E}=-\ln U$, with $U$ uniform on $(0,1)$, gives the familiar Gillespie waiting time. When $\Theta$ changes after a local event, we keep the same threshold and continue accumulating with the new rate. We draw a new macro threshold only after a macro event. Exact calculation of its first crossing defines the ideal process; the time-grid calculation below approximates that crossing. Section~\ref{sec:surrogate} shows that the ideal process is exact for the surrogate network with uniform inter-community connections.

The scheduler divides time into intervals of length $\tau$. On each interval, it approximates the varying macro rate by its value at the midpoint. Let $A$ be the amount already accumulated since the previous macro event. For the current interval $[t,t_g)$, copy every local simulator, including its random-number state, and advance each copy to the midpoint. The copies give an estimate $\Theta_{\mathrm{mid}}$ of the total macro rate. The estimated increase over the remaining interval is $\Theta_{\mathrm{mid}}(t_g-t)$. If $A$ plus this increase is below $\michael{\mathcal E}{E}$ (the random macro threshold), the real simulators move to $t_g$. Otherwise, estimate the event time by
\begin{equation}
\label{eq:interpolation}
t^\ast=t+\frac{\michael{\mathcal E}{E}-A}{\Theta_{\mathrm{mid}}}.
\end{equation}
Here, $\michael{\mathcal E}{E}$ is the random macro threshold and $A$ is the amount already accumulated toward it. The numerator $\michael{\mathcal E}{E}-A$ is what remains; dividing by the estimated rate gives the remaining waiting time. For example, if $\michael{\mathcal E}{E}=1$, $A=0.4$, and $\Theta_{\mathrm{mid}}=0.2$, the estimated wait is three time units. The scheduler advances the real simulators to $t^\ast$, selects the source and receiving communities using their states at that time, chooses the susceptible target, and imports the infection. Algorithm~\ref{alg:scheduler} gives the full sequence.

\begin{algorithm}[tbp]
\color{black}
\caption{\textcolor{black}{Grid-synchronised two-scale simulation of one realisation.}}
\label{alg:scheduler}
\small

\textbf{Input:} local graphs $G_i$, weights $W_{ij}$, rates, step $\tau$, horizon $T_{\mathrm{end}}$, and initial state.
\begin{enumerate}
\item Set $t=0$ and accumulated rate $A=0$; draw $\michael{\mathcal E}{E}\sim\Exp(1)$.
\item While $t<T_{\mathrm{end}}$ and infection can still spread, set $t_g=\min(t+\tau,T_{\mathrm{end}})$ and process the interval up to $t_g$.
\item Copy every local simulator, including its random-number state. Advance the copies to the midpoint of the remaining interval and calculate $\Theta_{\mathrm{mid}}$ from~\eqref{eq:macro_rate}.
\item If $A+\Theta_{\mathrm{mid}}(t_g-t)<\michael{\mathcal E}{E}$, advance the real simulators to $t_g$, add this amount to $A$, set $t=t_g$, and return to step 2.
\item Otherwise, calculate $t^\ast$ from~\eqref{eq:interpolation} and advance the real simulators to $t^\ast$.
\item Select $(i,j)$ with probability $\lambda_{ij}(t^\ast)/\Theta(t^\ast)$. Select $v$ uniformly from $S_j(t^\ast)$ and import infection at $v$. Update the receiving local simulator, keeping its unused exponential threshold and selecting the local event from the new rates.
\item Set $t=t^\ast$ and $A=0$, draw a new $\michael{\mathcal E}{E}\sim\Exp(1)$, and process the rest of the current interval in the same way.
\end{enumerate}

\end{algorithm}

The copied random-number states are important. Unless an import occurs first, a copied simulator follows the same path that the real simulator will follow to the midpoint. The midpoint estimate is exact on an interval with no local event. Its timing error and its look-ahead of at most $\tau/2$ decrease to zero as $\tau$ decreases.

The time grid is needed because the scheduler treats local simulators as black boxes. If they also report their next local-event times, the scheduler can merge those times and accumulate the piecewise-constant value of $\Theta$ exactly. That event-driven version removes the numerical synchronisation approximation. The experiments in this article use the implemented grid version with $\tau=1$.

\subsection{Computational structure and intended use}
\label{sec:cost}

\Full{} stores all local and cross-community channels, and its work per event grows with node degree \cite{kuryliak2025simulating}. \TS{} stores the local graphs and an $M\times M$ matrix of aggregate weights. The conference paper discusses this reduction in the information stored for cross-community contacts \cite{kuryliak2025towards}. Its importance depends on the problem. Replacing many individual cross-community edges by one weight can reduce storage. If only a few such edges exist, this saving may be small, and the local graphs still have to be stored.

The separate local simulators also provide an opportunity for parallel execution. Once the scheduler has chosen the next target time ($t_g$ or $t^\ast$), the local simulators can advance independently to that same time. The midpoint copies can also be processed in parallel. The macro layer must then collect the counts and apply any scheduled import before continuing. The implementation studied here advances the local simulators sequentially; parallel execution within one run is a possible extension, not a measured speedup. Local simulators can also use different community-specific settings.

The grid scheduler adds work by copying local states and coordinating the communities. Therefore, neither lower memory use nor faster execution is guaranteed. The method is useful when individual cross-community contacts are unavailable or too numerous to represent explicitly.

% English translation of the current Ukrainian main, active Sections 1--4.
% Translation of active Section 4 in the current Ukrainian main.
\section{Theoretical analysis of the two-scale method}
\label{sec:analysis}

The previous section described the two-scale simulation algorithm. The remaining question is how closely its results represent infection spread in the original network. Although the local simulators reproduce the within-community processes exactly, the macro layer uses aggregate connections instead of individual inter-community edges.

In this section, we first identify the network that mathematically corresponds to the ideal two-scale process. This establishes what changes when contacts are aggregated and separates these changes from the numerical synchronisation error. We then investigate how losing information about bridge nodes affects the first transmission time and introduce a characteristic for analysing accuracy in the experiments.

\subsection{Equivalence to simulation on a network with uniform inter-community connections}
\label{sec:surrogate}

The two-scale method replaces explicit inter-community contacts with aggregate infection rates while retaining exact stochastic dynamics within communities. This raises the question of whether the resulting process can be represented as exact simulation on a contact network and, if so, what that network looks like.

To answer this question, consider the ideal two-scale process, in which the time of each macro event is determined exactly by~\eqref{eq:hazard_integral}, without the time-grid approximation. This process is a continuous-time Markov chain (CTMC), because its transition rates depend on the current system state. Local transitions have the same rates as in the full-network model, while inter-community infections are determined by the aggregate rates~\eqref{eq:macro_rate} and the recipient-node selection rule.

We show that, when the susceptible recipient node is selected uniformly, this process is equivalent to exact simulation on a surrogate network with uniform inter-community connections, $\tilde G$ (Figure~\ref{fig:schematic}c). In this network, the total weight of each inter-community connection is distributed equally among all node pairs in the corresponding communities.

\begin{proposition}[Equivalence for a network with uniform inter-community connections]
\label{prop:surrogate}
Let $\tilde G$ have the same within-community contacts and their rates as the original network $G$. Replace each inter-community bundle $E_{ij}$ with the full set of channels $V_i\times V_j$, assigning each channel $(u,v)$ the rate
\[
\tilde\beta_{uv}
=
\frac{\beta_{\mathrm{out}}W_{ij}}{N_iN_j},
\qquad u\in V_i,\;v\in V_j,\;i\neq j.
\]
If the simulator selects the import target uniformly among susceptible nodes, the ideal two-scale process has the same rates for all transitions as the full-network process on $\tilde G$. Thus, from the same initial state, its distribution coincides with that of full-network Gillespie on $\tilde G$.
\end{proposition}

\begin{proof}
Local infections and removals have the same rates in both processes. It therefore suffices to compare the inter-community infection rates.

Consider any susceptible node $v\in V_j$ in system state $x$. In the two-scale process, imports from community $i$ to $j$ occur at rate $\lambda_{ij}(x)$, and the probability of selecting node $v$ is $1/s_j(x)$. Its total infection rate from other communities is therefore
\[
\begin{aligned}
\sum_{i\neq j}\lambda_{ij}(x)\frac{1}{s_j(x)}
&=
\sum_{i\neq j}
\beta_{\mathrm{out}}W_{ij}
\frac{y_i(x)}{N_i}
\frac{s_j(x)}{N_j}
\frac{1}{s_j(x)}
\\
&=
\sum_{i\neq j}
\sum_{u\in V_i}
\frac{\beta_{\mathrm{out}}W_{ij}}{N_iN_j}
\ind\{x_u=I\}.
\end{aligned}
\]

The last expression is the total infection rate of node $v$ through all inter-community channels of the surrogate network $\tilde G$ that lead to it from infected nodes. If $s_j(x)=0$, the community has no susceptible nodes, and inter-community infection is impossible in both models.

Thus, the rates of all possible transitions coincide. The generators of the two CTMCs are therefore identical, giving the same process distribution from the same initial state.
\end{proof}

\begin{corollary}
\label{cor:exact}
If every node in community $i$ is already connected to every node in community $j$ in the original network, and all these channels have the same rate for a given community pair, then $\tilde G=G$. In this case, the ideal two-scale method is also exact for the original network.
\end{corollary}

The ideal two-scale method therefore has an exact network interpretation: it reproduces the process on $\tilde G$. In general, however, the original and surrogate networks differ. The established equivalence does not by itself guarantee similar simulation results on these networks. To identify possible reasons for discrepancies, we examine which properties of inter-community contacts are preserved by their replacement and which change.

\subsection{Mechanisms of discrepancy from the full-network model}
\label{sec:sources}

By construction, the surrogate network $\tilde G$ preserves within-community contacts and the total coupling rate between each community pair:
\[
\sum_{(u,v)\in V_i\times V_j}\tilde\beta_{uv}
=
\beta_{\mathrm{out}}W_{ij}.
\]

This equality concerns the total rate of all channels, however, rather than only those through which infection can be transmitted in the current system state. The actual infection rate depends on the endpoint states of each contact. Equal total inter-community weights therefore do not guarantee the same infection dynamics.

Replacing real inter-community contacts by averaged connections changes three things: the transmission rate, the recipient node, and the possible number of imports. A fourth source of error is the approximate calculation of event times in the implementation.

\paragraph{1. The states of the real bridge nodes are not used.}
In the full network, the transmission rate from $V_i$ to $V_j$ is
\begin{equation}
\label{eq:exact_rate}
a^\star_{ij}(x)=\beta_{\mathrm{out}}\sum_{(u,v)\in E_{ij}}w_{uv}\,\ind\{x_u=I,\,x_v=S\},
\end{equation}
Only contacts from infected nodes to susceptible nodes contribute to this sum. If all bridge nodes in the source are still susceptible, transmission is impossible. In the two-scale method, the rate depends on the overall infected fraction, so infection of any source node can contribute to export. The method can transmit infection before any real bridge node is infected. We refer to this as \emph{bridge bypass} (Section~\ref{sec:single_export}).

The aggregate rate does not necessarily exceed the actual rate. If infected nodes are concentrated at the inter-community interface, the opposite can occur. Comparing rates in an individual state is therefore insufficient to establish the overall direction of a shift in transmission time.

\paragraph{2. Import location.}
In the full network, infection reaches the endpoint of the contact through which it was transmitted. The two-scale method selects any susceptible node with equal probability.
As a result, infection may reach a node that has no incoming contact from the corresponding source community in the original network. The location of the first infection affects subsequent spread within the community, including the time taken to reach bridge nodes leading to further communities. Changing the import-placement rule can therefore affect both local dynamics and infection spread through a sequence of communities.

\paragraph{3. Additional inter-community infections.}
In the SI model, each node can move from $S$ to $I$ only once. Each directed channel in the original network can therefore cause at most one infection of its target node. Consequently, at most $|E_{ij}|$ such infections can occur through the bundle $E_{ij}$.

The two-scale method does not track individual inter-community channels. After the first import, the rate $\lambda_{ij}$ can remain positive as long as the source community has infected nodes and the recipient community has susceptible nodes.

For example, if the communities are connected by just one directed channel, it can cause at most one infection in the full network. In contrast, \TS{} can import infection into several different nodes of the recipient community. These additional infections can accelerate internal spread and subsequent transmission to other communities.

\paragraph{4. Numerical synchronisation error.}
As described in Section~\ref{sec:timing}, the scheduler estimates the rate from the system state at the midpoint of a synchronisation interval and uses it to approximate the integral~\eqref{eq:hazard_integral}. The implemented \TS{} can therefore differ even from the ideal two-scale process on $\tilde G$.

% Ukrainian editorial note retained for review: is this remark needed?
\begin{remark}[The first import with randomly placed contacts]
\label{rem:placement}
Suppose that the endpoints of inter-community channels are selected uniformly and independently of the internal graphs and other contact sets, as in our tests. If the recipient community is still fully susceptible, the target of its first import is also uniform in the full network, after averaging over random contact placement. The two-scale method has the same distribution of the first target. For an individual fixed network, the targets need not coincide. The same distribution of the first target does not imply the same distribution of the first import time. Our tests show small errors in a number of configurations where the spread of internal infection times is small compared with the mean transmission time through one inter-community channel (Section~\ref{sec:res_descriptor}).
\end{remark}

Replacing inter-community contacts thus changes not only the nodes through which infection can be transmitted, but also the possible locations and number of infections. These mechanisms can act simultaneously and should not be treated as independent additive components of the total error. Analysing them alone does not establish either the overall direction or the magnitude of discrepancies between the models.

\subsection{Effect of contact aggregation on systematic shifts in export time}
\label{sec:single_export}

We now examine how replacing individual inter-community contacts with an aggregate connection affects the first transmission time to another community. In particular, we ask whether aggregation can systematically shift the first export time even when the mean transmission rate is preserved.

To isolate the effect of losing information about bridge nodes, consider a simplified network: a connected source community $V_0$ with $n$ nodes and one external susceptible node $\ell$, called a leaf. The leaf is connected to a set $B\subset V_0$ of $k$ bridge nodes, with each channel having rate $\beta_{\mathrm{out}}$. The source follows SI dynamics, and the leaf does not affect internal infection spread before it becomes infected.

In the full-network model, the leaf can become infected only after at least one real bridge node becomes infected. In \TS{}, the aggregate coupling weight is $W_{0\ell}=k$, so before the leaf is infected, the export rate is
\[
\lambda_{0\ell}(t)
=
\beta_{\mathrm{out}}k\ibar_0(t).
\]
The possibility of export in the two-scale model therefore depends on the overall number of infected source nodes, rather than the states of the individual bridges.

Consider the case in which the $k$ bridge nodes are selected uniformly and independently of the internal dynamics. This allows us to compare first-export-time distributions while accounting for random placement of inter-community contacts.

\begin{proposition}[Stochastic ordering of the first export time]
\label{prop:ordering}
Let the bridge set $B$ be chosen uniformly from all $k$-element subsets of $V_0$, independently of the internal graph, initial state, and event randomness. Let $T^{\Full}$ and $T^{\TS}$ be the first infection times of the leaf in the full-network and ideal two-scale models, respectively. Then
\[
\mathbb{P}(T^{\Full}>t)
\geq
\mathbb{P}(T^{\TS}>t),
\qquad t\geq 0.
\]

Thus, with randomly placed bridges, the first export time in \TS{} is stochastically no greater than the corresponding time in \Full{}. The same ordering holds for its expectation and quantiles.
\end{proposition}

\begin{proof}
Before the leaf becomes infected, infection spread within $V_0$ is independent of external transmission. Both models can therefore be compared using the same internal SI trajectory $X_0$ of the source community.

Let $\tau_u$ be the infection time of node $u$, and let $A_u(t)=\max\{t-\tau_u,0\}$ be the time it has spent infected by time $t$.

In the full-network model, the cumulative export rate up to time $t$ is determined only by infected bridge nodes:
\[
H_B(t)
=
\beta_{\mathrm{out}}\sum_{u\in B}A_u(t).
\]

In \TS{}, the corresponding cumulative rate is determined by all source nodes:
\[
\bar H(t)
=
\beta_{\mathrm{out}}\frac{k}{n}
\sum_{u\in V_0}A_u(t).
\]

With a uniform, independent choice of $B$, each source node has probability $k/n$ of belonging to the bridge set. Conditional on the trajectory $X_0$, therefore,
\[
\mathbb{E}_B[H_B(t)\mid X_0]
=
\bar H(t).
\]

The cumulative rate in \TS{} is thus the expected cumulative rate in \Full{} over random bridge placement. Equality of these quantities does not, however, imply equality of the export-time distributions.

The probability of no export by time $t$ is the exponential of the negative cumulative rate. Since $z\mapsto e^{-z}$ is convex, Jensen's inequality gives
\[
\begin{aligned}
\mathbb{P}(T^{\Full}>t\mid X_0)
&=
\mathbb{E}_B[e^{-H_B(t)}\mid X_0]
\\
&\geq
e^{-\mathbb{E}_B[H_B(t)\mid X_0]}
\\
&=
e^{-\bar H(t)}
\\
&=
\mathbb{P}(T^{\TS}>t\mid X_0).
\end{aligned}
\]

Because this inequality holds for every internal trajectory $X_0$, it also holds after averaging over all possible trajectories, proving the proposition.
\end{proof}

This result shows that the aggregate rate matches the expected actual rate, yet the transmission-time distributions can differ systematically. The reason is the nonlinear dependence of the probability of no event on the cumulative rate: averaging rates is not equivalent to averaging the corresponding probabilities.

In addition to shifting time, aggregation can violate constraints on admissible transmission routes. Let $T_B=\min_{u\in B}\tau_u$ be the infection time of the first real bridge node. In the full-network model, the leaf cannot become infected before $T_B$. In ideal \TS{}, however, the probability of such an early export is
\begin{equation}
\label{eq:bypass}
\mathbb{P}(T^{\TS}<T_B)
=
\mathbb{E}\!\left[
1-\exp\left(
-\beta_{\mathrm{out}}\frac{k}{n}
\int_0^{T_B}y_0(s)\,ds
\right)
\right].
\end{equation}

This expression gives the probability that the aggregate rate causes an export before any real bridge node becomes infected. It shows that even when export-time distributions are close, the two models can differ in whether their transmission paths are admissible.

For randomly placed bridges, we have therefore established the direction of the stochastic shift in the first export time, but not its magnitude. The next question is which properties of internal spread make this discrepancy small or substantial.

% The former Section 4 temporal-descriptor drafts are commented out in the
% Ukrainian source; the active descriptor is now in Section 5 below.

\section{Methodology for experimental evaluation}
\label{sec:design}

\subsection{Experimental networks and parameters}

The theoretical analysis in the previous section showed how aggregating inter-community contacts changes infection transmission and under which conditions a stochastic shift in the first export time arises. However, it does not determine the magnitude of the discrepancies between the models or their behaviour after successive inter-community crossings.

To investigate these questions, we compare \Full{} and \TS{} in four experimental series. We first consider transmission to a single external node and then infection spread along a chain of communities. In each setting, we examine how internal network structure affects the discrepancies between the models.

All experiments use SI dynamics, unit contact weights, $\beta_{\mathrm{in}}=\beta_{\mathrm{out}}=\beta$, and a synchronisation step of $\tau=1$ (Algorithm~\ref{alg:scheduler}).

\paragraph{Clique--leaf.}
To compare the first export times of the two methods with minimal influence from internal network structure, we use a complete graph. The receiving community need not be represented in full: for this purpose, it is sufficient to represent export through the bridge nodes and a single external leaf.

We attach one external node (the leaf) to a complete graph (clique) with $n=100$ nodes through $k=2$ distinct bridge nodes in the clique. We perform 10\,000 runs of each model with $\beta=0.005$ and a time horizon of 1000. In this and all other experiments, bridge nodes are selected uniformly without replacement.

\paragraph{Random community--leaf.}
We next investigate how the internal structure of the source affects infection transmission to the leaf. We replace the clique with a connected random graph of $n=1000$ nodes and attach the leaf through $k\in\{1,5\}$ distinct bridge nodes.

We characterise internal structure by the average shortest-path length (ASPL). To generate the graph, we first construct a uniformly random tree and then add random edges to obtain the target ASPL within a tolerance of 0.1. We vary ASPL from 1 to 6 in steps of 0.25; ASPL 1 corresponds to a complete graph.

As ASPL increases, the mean degree of the generated graphs also decreases. This experiment therefore investigates the effects of internal path length and network sparsity.

The initially infected node is selected uniformly from all source nodes. For each configuration, we perform 400 runs of each model: 40 realisations of the internal graph with 10 replicates of inter-community contacts and random events. We use $\beta=0.003$ and a horizon of 30\,000.

\paragraph{Chain of cliques.}
We assess how discrepancies change after successive crossings in a chain of $L=10$ complete graphs, each with $n=1000$ nodes. Neighbouring communities are connected by one edge ($k=1$). Infection starts in community 0.

As in the first test, complete graphs provide rapid internal spread. This allows us to assess the effect of successive inter-community crossings in a baseline setting with minimal internal path lengths.

We use $\beta=0.003$, 400 runs of each model, and a horizon of 30\,000.

\paragraph{Chains of random communities.}
Finally, we combine successive inter-community crossings with variation in internal community structure. This series examines how community size and internal structure affect discrepancies between \Full{} and \TS{}, and how these discrepancies change along the chain.

We connect $L=10$ connected random communities in a line, with $k\in\{1,5\}$ edges between neighbouring communities. Infection starts in community 0.

We consider sizes $n\in\{50,100,250,1000\}$. Internal graphs are generated by the procedure described above, varying ASPL from 1 to 6 in steps of 1.

For each interface, we select $k$ nodes uniformly without replacement from each community and connect them in pairs. A node can participate in interfaces with both neighbouring communities.

For each configuration, we perform 400 runs of each model with $\beta=0.003$ and a horizon of 30\,000. \yulian{Two \Full{} runs ($n=50$, ASPL 6, $k=1$) were extended to 40\,000 using the same graphs, initial nodes and event seeds, to observe infection of the last community. Arrival-time summaries use the observed event times; unobserved arrivals were not replaced by the horizon.}{Two \Full{} runs were extended to 40\,000 to observe infection of the last community (Appendix~\ref*{app:repro}).}

\begin{table*}[!t]
\centering
\caption{Parameters of the experimental series. ASPL is the target average shortest-path length within a community.}
\label{tab:design}
\small
\begin{tabular}{
>{\raggedright\arraybackslash}p{0.17\linewidth}
>{\raggedright\arraybackslash}p{0.36\linewidth}
>{\raggedright\arraybackslash}p{0.06\linewidth}
>{\raggedright\arraybackslash}p{0.08\linewidth}
>{\raggedright\arraybackslash}p{0.20\linewidth}}
\toprule
Series & Structure & $k$ & $\beta$
& Runs per cell and model \\
\midrule
Clique--leaf
& Clique $n=100$ and one leaf
& 2 & 0.005 & 10\,000 \\

Random community--leaf
& $n=1000$, ASPL 1--6 in steps of 0.25,
one leaf
& 1, 5 & 0.003 & 400 (40 $\times$ 10) \\

Chain of cliques
& $L=10$, $n=1000$
& 1 & 0.003 & 400 (40 $\times$ 10) \\

Chains of random communities
& $L=10$, $n\in\{50,100,250,1000\}$,
ASPL 1--6 in steps of 1
& 1, 5 & 0.003 & 400 (40 $\times$ 10) \\
\bottomrule
\end{tabular}
\end{table*}

\michael{Each pair of \Full{} and \TS{} runs uses the same internal graphs, inter-community contacts, and initially infected node, but independent streams of random events. The chain-of-cliques benchmark is the ASPL-1, $n=1000$, $k=1$ configuration of the 48-cell chain sweep, rather than an additional independent configuration.}{Each pair of \Full{} and \TS{} runs uses the same internal graphs, inter-community contacts, and initially infected node, but independent streams of random events.} Inter-community contacts are generated independently for different interfaces and replicates.

\subsection{Metrics for comparing the models}

To assess discrepancies between \Full{} and \TS{}, we compare infection propagation times and cases of transmission before the bridge nodes become infected.

In the leaf tests, we measure its first infection time $T_{\mathrm{export}}$. For chains, let $T_c$ be the first infection time of community $c$, where $c=1,\ldots,9$. Since the chain has 10 communities, the arrival time in the last community is denoted by $T_9$.

Because infection spread is stochastic, the time of the same event varies between runs. The first infected nodes and infection paths also differ between runs. Comparing individual realisations is therefore insufficient to assess discrepancies between the models.

For each configuration, we perform repeated runs of \Full{} and \TS{}. Each run yields one time for the corresponding event: infection of the leaf or the first infection of a particular community. From these values, we form separate empirical distributions for the two models and compare them using three metrics.

\paragraph{Relative median shift.}
To determine how much earlier or later \TS{} transmits infection than \Full{}, we calculate the relative difference between the medians of the corresponding distributions:
\[
b_c =
100\,
\frac{
\operatorname{med}(T_c^{\TS})-
\operatorname{med}(T_c^{\Full})
}{
\operatorname{med}(T_c^{\Full})
}.
\]
A negative value means earlier infection arrival in \TS{}, while a positive value means later arrival. The median shift describes the direction and magnitude of the difference in typical times, but does not describe the whole distribution.

\paragraph{Normalised Wasserstein distance.}
The median shift captures only the difference between the medians and does not account for the rest of the distribution. To assess differences between empirical distributions of event times, we therefore use the Wasserstein-1 distance.

We sort the $R$ observed times of the corresponding event in ascending order separately for each model and calculate
\[
W_1 =
\frac{1}{R}
\sum_{r=1}^{R}
\left|
T_{(r)}^{\TS}-T_{(r)}^{\Full}
\right|,
\]
where $T_{(r)}^{\TS}$ and $T_{(r)}^{\Full}$ are the $r$th values in the ordered \TS{} and \Full{} samples, respectively.

To compare results with different characteristic propagation times, we normalise the distance by the mean time in \Full{}:
\[
\eW =
100\,\frac{W_1}{\overline{T}^{\Full}}.
\]
The resulting measure describes the difference between empirical distributions independently of the direction of the shift. The formula compares ordered samples, not individual pairs of runs.

In the leaf tests, both metrics are calculated for $T_{\mathrm{export}}$; in chains, they are calculated for the arrival times $T_c$.

\paragraph{Fraction of premature exports}
Similar infection-time distributions do not guarantee that an admissible sequence of infection transmission is preserved. In the full network, infection cannot reach a neighbouring community until at least one real bridge node in the source becomes infected. In contrast, \TS{} can transmit infection before any of these nodes becomes infected.

To assess how often this occurs, for each inter-community crossing we consider the set $B$ of real bridge nodes in the source community. Let $\tau_u^{\TS}$ be the infection time of node $u$ in the corresponding \TS{} run. The infection time of the first bridge is then
\[
T_B^{\TS}=\min_{u\in B}\tau_u^{\TS}.
\]

Let $T_{\mathrm{arr}}^{\TS}$ be the first infection time of the recipient community or leaf. We record a premature export if
\[
T_{\mathrm{arr}}^{\TS}<T_B^{\TS}.
\]

The fraction of premature exports is the percentage of such cases among all examined inter-community crossings in \TS{}. This metric describes how often a necessary condition for transmission through real bridge nodes is violated.

\paragraph{Discrepancies along the chain.}
For chains, we additionally compare the relative shifts in the mean durations of the first crossing, $T_1$, and the subsequent crossings, $T_9-T_1$. The first crossing allows us to assess discrepancies associated with export from the initial community, while later crossings can also reflect the effects of import placement and surplus imports. This comparison tests whether discrepancies become stronger after the first inter-community crossing.

\subsection{A measure of the ratio of time scales}
\label{sec:descriptor}

Discrepancies between \Full{} and \TS{} may depend not only on network structure, but also on the relative speeds of infection spread within and between communities. In particular, aggregation does not account for when individual bridge nodes become infected, which can matter when internal spread is slow and uneven.

We characterise temporal heterogeneity of infection by the dimensionless measure
\begin{equation}
\label{eq:rho}
\rho=\beta_{\mathrm{out}}\sigma_T,
\end{equation}
where $\sigma_T$ is the standard deviation of node infection times within a community. This measure describes the spread of internal infection times relative to the mean transmission time through one continuously active inter-community channel, $1/\beta_{\mathrm{out}}$.

For each configuration, we calculate $\sigma_T$ from node infection times in each source community of the \Full{} runs, then average these values over communities and runs.

We examine how $\rho$ relates to the median shift and normalised Wasserstein distance for the arrival time in the last community, $T_9$. In particular, we compare how well $\rho$ and ASPL explain discrepancies between the models across community sizes, structures, and numbers of inter-community edges.

A small $\rho$ means that the spread of internal infection times is small relative to the transmission time through one channel.

The measure $\rho$ is not a direct error metric. We consider it a possible explanatory characteristic of the discrepancies. Since $\sigma_T$ is calculated from \Full{}, we do not test whether error can be predicted from local simulations alone.

% User's Ukrainian editorial note: not sure whether this subsection is needed.
\subsection{Statistical evaluation}

Because the simulations are stochastic and the number of runs is limited, the estimated discrepancies between \Full{} and \TS{} have sampling uncertainty. We use a block bootstrap to assess this uncertainty and a permutation test to test differences between infection-time distributions.

\paragraph{Block bootstrap.}
Replicates on the same graph share its structure and are therefore not independent realisations of a network configuration. To account for this dependence, the bootstrap samples whole network blocks with replacement, keeping all ten replicates and their corresponding pairs of \Full{} and \TS{} runs together.

For chain summary statistics, we form 2000 bootstrap samples and recalculate the measures of interest. The limits of the 95\% confidence intervals are the 2.5th and 97.5th percentiles of the resulting estimates. This assesses uncertainty in the magnitude and direction of discrepancies between the models.

\paragraph{Permutation test.}
The Wasserstein distance can be positive even for two samples from the same distribution, because each contains a finite number of random observations. We therefore need to test whether the observed distance exceeds the level that can arise from sampling variability.

For each pair of runs, we randomly swap the observed \Full{} and \TS{} infection times with probability $1/2$. We recalculate the normalised Wasserstein distance from the resulting samples. Repeating this procedure 2000 times gives a distribution of distances under random reassignment of results between models.

Under the null hypothesis that \Full{} and \TS{} results are exchangeable within pairs, we use this distribution to assess the statistical significance of the observed distance. Its 95th percentile defines the permutation threshold, and the fraction of permutations with a distance at least as large as the observed one estimates the $p$-value.

\michael{Exceeding the threshold indicates a statistically significant discrepancy under the assumed null model. Failure to reject this null does not establish equivalence of the two distributions.}{Exceeding the threshold indicates a statistically significant discrepancy under the assumed null model.}

% User's Ukrainian editorial note: not sure whether this subsection is needed.
\subsection{Validation of the Full implementation}

To validate the full-network simulator, we compare the simulated mean infection time of the first bridge node, $T_B$, with the analytical value for a clique (Appendix~\ref*{app:clique}). This check uses four clique sizes and two numbers of inter-community edges.

We also check the mean transmission time to the leaf after the single bridge becomes infected. For one continuously active channel, its theoretical value is $1/\beta$.

These checks assess the agreement of the \Full{} implementation with known analytical results, but do not establish the accuracy of the two-scale approximation.

\section{Results}
\label{sec:results}

In this section, we compare \Full{} and \TS{} in terms of the timing of infection spread and the admissibility of inter-community transmission. We use the experimental series described in Section~\ref{sec:design}. We first consider the first export of infection to determine the consequences of losing information about bridge nodes. We then examine spread along a chain of communities, the effects of network characteristics, and the relationship between the discrepancies and the ratio of time scales.

The chain comparisons cover 48 configurations, with 400 paired \Full{} and \TS{} runs in each, for a total of 19\,200 pairs.

\subsection{Temporal accuracy and admissibility of inter-community transmission}
\label{sec:res_clique}

The first experiment tests whether \TS{} can reproduce the infection export time accurately without retaining information about the specific bridge nodes. We consider a clique--leaf network, in which internal spread is rapid.

In this case, the export-time distributions were close (Figure~\ref{fig:clique_leaf}). The mean time to infection of the leaf was 110.18 in \Full{} and 112.62 in \TS{}, and the normalised Wasserstein distance was $\eW=2.95\%$. The paired difference in means was $+2.43$, with a 95\% bootstrap interval of $[-0.35,5.27]$, which contains zero. Thus, this experiment did not detect a significant systematic shift in the mean export time.

\begin{figure*}[!tbp]
\centering
% BEGIN EDITABLE FIGURE clique_leaf
% Editable PGFPlots version of clique_leaf_endpoint_and_mechanism.
\begin{tikzpicture}
\begin{groupplot}[
  group style={group size=3 by 1, horizontal sep=1.55cm},
  set layers=standard,
  paper axis, width=.24\linewidth, height=3.7cm, scale only axis,
  legend style={font=\scriptsize, draw=none},
]
\nextgroupplot[
  title={\textbf{a} Export times},
  xlabel={total export time}, ylabel={empirical CDF},
  xmin=0, xmax=600, ymin=0, ymax=1.01,
  legend style={at={(.98,.36)}, anchor=south east},
  extra description/.code={
    \DrawExactECDF{figures_pgfplots/leaf_clique_export_full_path.dat}{FullBlue}{\pgfsetdash{}{0pt}}
    \DrawExactECDF{figures_pgfplots/leaf_clique_export_twoscale_path.dat}{ScaleOrange}{\pgfsetdash{{3pt}{3pt}}{0pt}}
  },
]
  \addlegendimage{color=FullBlue, solid, no markers, line width=.8pt}
  \addlegendentry{Full}
  \addlegendimage{color=ScaleOrange, dashed, no markers, line width=.8pt}
  \addlegendentry{TwoScale}
  \node[anchor=south east, align=right, fill=white, draw=black!20, inner sep=2pt, font=\scriptsize]
    at (rel axis cs:.97,.07) {$W_1=3.249$\\$D_{\rm KS}=0.0142$};
\nextgroupplot[
  title={\textbf{b} Export minus bridge},
  xlabel={$T_{\mathrm{export}}-T_B$}, ylabel={empirical CDF},
  xmin=-20, xmax=100, ymin=0, ymax=.55,
  extra description/.code={
    \DrawExactECDF{figures_pgfplots/leaf_clique_signed_delay_full_path.dat}{FullBlue}{\pgfsetdash{}{0pt}}
    \DrawExactECDF{figures_pgfplots/leaf_clique_signed_delay_twoscale_path.dat}{ScaleOrange}{\pgfsetdash{{3pt}{3pt}}{0pt}}
  },
]
  \path[fill=black!8] (axis cs:-20,0) rectangle (axis cs:0,.55);
  \addplot+[black!75, no markers, line width=.35pt] coordinates {(0,0) (0,.55)};

  \node[anchor=north west, font=\scriptsize, text=black!65] at (axis cs:-18,.51) {export before bridge};
\nextgroupplot[
  title={\textbf{c} Bridge bypass},
  ylabel={bridge-bypass runs (\%)},
  xmin=.5, xmax=2.5, ymin=-.03, ymax=1.35, clip=false,
  xtick={1,2}, xticklabels={Full,TwoScale},
]
  \addplot+[
    only marks, mark=*, mark size=1.7pt, color=black!80,
    error bars/.cd, y dir=both, y explicit,
  ] table[x=x,y=rate_pct,
      y error minus expr=\thisrow{rate_pct}-\thisrow{low_pct},
      y error plus expr=\thisrow{high_pct}-\thisrow{rate_pct}]
      {figures_pgfplots/leaf_clique_bypass.dat};
  \node[anchor=south, align=center, font=\scriptsize] at (axis cs:1,.14) {0/10000\\(0.00\%)};
  \node[anchor=south, align=center, font=\scriptsize] at (axis cs:2,1.18) {98/10000\\(0.98\%)};
\end{groupplot}
\end{tikzpicture}
% END EDITABLE FIGURE clique_leaf
\caption{%
First export of infection in the clique--leaf network ($n=100$, $k=2$, $\beta=0.005$, 10\,000 runs of each model).
(a) Empirical distribution functions of the export time.
(b) Distributions of the difference $T_{\mathrm{export}}-T_B$; negative values correspond to export before either real bridge node becomes infected.
(c) Fraction of runs with this violation of the transmission order and 95\% confidence intervals.
}
\label{fig:clique_leaf}
\end{figure*}

The similarity of the time distributions in this experiment is consistent with rapid and uniform infection spread within the clique. Because of its symmetric structure, infection spreads quickly among the nodes regardless of their positions, including the bridge nodes.

The expected time to infection of the first bridge is only $\mathbb{E}[T_B]=8.36$, whereas the mean export time exceeds 110. Thus, the actual and aggregate transmission rates differ for a relatively short period compared with the overall waiting time for export. This is consistent with the small discrepancy observed between the models.

However, similar time distributions do not imply that the admissible infection route is preserved. In 98 of the 10\,000 \TS{} runs (0.98\%), the leaf was infected before either of the two real bridge nodes. This event order is impossible in \Full{}. The aggregate interface can reproduce the timing of transmission accurately while allowing transitions that are impossible in the original network.

To determine how this effect depends on internal structure, we replace the clique with connected random graphs having different average shortest path lengths (ASPL). In this series, graph density also decreases as ASPL increases.

In dense networks with short paths, premature exports remained rare: their fraction did not exceed 0.5\% for ASPL $\leq 1.75$. It increased with ASPL and, at ASPL 6, reached 32.5\% for both values of $k$ examined (Figure~\ref{fig:source_sweep}).

\begin{figure*}[!tbp]
\centering
% BEGIN EDITABLE FIGURE source_sweep
% Editable PGFPlots version of fig_aspl_export_fidelity.
\begin{tikzpicture}
\begin{groupplot}[
  group style={group size=3 by 2, horizontal sep=1.55cm, vertical sep=1.35cm},
  paper axis, width=.24\linewidth, height=3.7cm, scale only axis,
  xmin=.85, xmax=6.15, xtick={1,2,...,6},
]
\nextgroupplot[
  title={$k=1$: export time}, ylabel={mean export time},
  xticklabels=\empty, ymin=0, ymax=1643.82646974,
  legend style={font=\scriptsize, draw=none, at={(.98,.98)}, anchor=north east},
]
  \addplot+[color=FullBlue, mark=*, mark size=1.1pt, mark options={fill=FullBlue,draw=FullBlue}, line width=.6pt,
    error bars/.cd, y dir=both, y explicit]
    table[x=aspl,y=full_mean,y error minus expr=\thisrow{full_mean}-\thisrow{full_low},y error plus expr=\thisrow{full_high}-\thisrow{full_mean}]
    {figures_pgfplots/leaf_source_sweep_k1.dat};
  \addlegendentry{Full}
  \addplot+[color=ScaleOrange, mark=square*, mark size=1.1pt, mark options={fill=ScaleOrange,draw=ScaleOrange}, line width=.6pt,
    error bars/.cd, y dir=both, y explicit]
    table[x=aspl,y=twoscale_mean,y error minus expr=\thisrow{twoscale_mean}-\thisrow{twoscale_low},y error plus expr=\thisrow{twoscale_high}-\thisrow{twoscale_mean}]
    {figures_pgfplots/leaf_source_sweep_k1.dat};
  \addlegendentry{TwoScale}
\nextgroupplot[title={$k=1$: $W_1$ error}, ylabel={$W_1/\overline T_{\rm Full}$ (\%)}, xticklabels=\empty, ymin=0, ymax=13.026367574]
  \addplot+[color=ScaleOrange, mark=*, mark size=1.1pt, mark options={fill=ScaleOrange,draw=ScaleOrange}, line width=.6pt]
    table[x=aspl,y=w1]{figures_pgfplots/leaf_source_sweep_k1.dat};
\nextgroupplot[title={$k=1$: bridge bypass}, ylabel={export before any bridge (\%)}, xticklabels=\empty, ymin=0, ymax=39.9]
  \addplot+[color=ScaleOrange, mark=*, mark size=1.1pt, mark options={fill=ScaleOrange,draw=ScaleOrange}, line width=.6pt]
    table[x=aspl,y=bypass]{figures_pgfplots/leaf_source_sweep_k1.dat};
  \addplot+[color=ScaleOrange, mark=none, line width=.45pt,
    error bars/.cd, y dir=both, y explicit]
    table[x=aspl,y=bypass,y error minus expr=\thisrow{bypass}-\thisrow{bypass_low},y error plus expr=\thisrow{bypass_high}-\thisrow{bypass}]
    {figures_pgfplots/leaf_source_sweep_k1_bypass_ci.dat};
\nextgroupplot[title={$k=5$: export time}, xlabel={target ASPL}, ylabel={mean export time}, ymin=0, ymax=1643.82646974]
  \addplot+[color=FullBlue, mark=*, mark size=1.1pt, mark options={fill=FullBlue,draw=FullBlue}, line width=.6pt,
    error bars/.cd, y dir=both, y explicit]
    table[x=aspl,y=full_mean,y error minus expr=\thisrow{full_mean}-\thisrow{full_low},y error plus expr=\thisrow{full_high}-\thisrow{full_mean}]
    {figures_pgfplots/leaf_source_sweep_k5.dat};
  \addplot+[color=ScaleOrange, mark=square*, mark size=1.1pt, mark options={fill=ScaleOrange,draw=ScaleOrange}, line width=.6pt,
    error bars/.cd, y dir=both, y explicit]
    table[x=aspl,y=twoscale_mean,y error minus expr=\thisrow{twoscale_mean}-\thisrow{twoscale_low},y error plus expr=\thisrow{twoscale_high}-\thisrow{twoscale_mean}]
    {figures_pgfplots/leaf_source_sweep_k5.dat};
\nextgroupplot[title={$k=5$: $W_1$ error}, xlabel={target ASPL}, ylabel={$W_1/\overline T_{\rm Full}$ (\%)}, ymin=0, ymax=13.026367574]
  \addplot+[color=ScaleOrange, mark=*, mark size=1.1pt, mark options={fill=ScaleOrange,draw=ScaleOrange}, line width=.6pt]
    table[x=aspl,y=w1]{figures_pgfplots/leaf_source_sweep_k5.dat};
\nextgroupplot[title={$k=5$: bridge bypass}, xlabel={target ASPL}, ylabel={export before any bridge (\%)}, ymin=0, ymax=39.9]
  \addplot+[color=ScaleOrange, mark=*, mark size=1.1pt, mark options={fill=ScaleOrange,draw=ScaleOrange}, line width=.6pt]
    table[x=aspl,y=bypass]{figures_pgfplots/leaf_source_sweep_k5.dat};
  \addplot+[color=ScaleOrange, mark=none, line width=.45pt,
    error bars/.cd, y dir=both, y explicit]
    table[x=aspl,y=bypass,y error minus expr=\thisrow{bypass}-\thisrow{bypass_low},y error plus expr=\thisrow{bypass_high}-\thisrow{bypass}]
    {figures_pgfplots/leaf_source_sweep_k5_bypass_ci.dat};
\end{groupplot}
\end{tikzpicture}
% END EDITABLE FIGURE source_sweep
\caption{%
Effect of the source's internal structure on the first export to the leaf ($n=1000$, $\beta=0.003$, $k=1$ in the top row and $k=5$ in the bottom row; 400 paired runs for each configuration).
Left: mean export time.
Centre: normalised Wasserstein distance $\eW$ between the export-time distributions.
Right: fraction of \TS{} exports occurring before infection of the first real bridge; such events are impossible in \Full{}.
Error bars show 95\% block-bootstrap intervals.
}
\label{fig:source_sweep}
\end{figure*}

The timing of export also changed. At ASPL 6, the mean export time in \TS{} was 8.3\% lower than in \Full{} for $k=1$ and 6.2\% lower for $k=5$. The corresponding 95\% bootstrap intervals were $[-13.1,-3.5]\%$ and $[-10.3,-2.3]\%$. \michael{Some configurations had positive empirical mean shifts; for example, the shift was $+5.2\%$ at ASPL 1 and $k=1$. The stochastic ordering in Proposition~\ref{prop:ordering} concerns ideal-process distributions, not finite-sample estimates or the grid-based implementation.}{A positive shift in the mean export time was not observed in any configuration in this series.}

The direction of this shift is consistent with Proposition~\ref{prop:ordering} for the ideal two-scale process: averaging rates over randomly selected bridge nodes can result in systematically earlier export.

\yulian{}{Sampling variability must be considered when interpreting the Wasserstein distance. \michael{For independent samples of size $R$ from an exponential distribution with mean $\mu$, the large-$R$ approximation is $\mathbb{E}[W_1]/\mu\approx\sqrt{\pi/R}$, giving about $8.9\%$ for $R=400$ (Appendix~\ref*{app:floor}). This is an illustrative sampling baseline, not the critical value for our paired permutation tests.}{For nearly exponential export times with $R=400$, the expected normalised distance between two independent samples from the same distribution is approximately $\sqrt{\pi/R}\approx9\%$ (Appendix~\ref*{app:floor}).} Small values of $\eW$ for a single export therefore do not, by themselves, indicate a difference between the models.}

Unlike the temporal distance, every premature export directly indicates a violation of the admissible transmission order. For example, at ASPL 4 and $k=1$, such violations occurred in 19.75\% of runs, although the time distributions remained relatively close.

The first export demonstrates two distinct consequences of aggregation. Even when the transmission times are close, the model can violate the physical constraints imposed by inter-community contacts. When internal spread becomes slower and less uniform, this loss of information can also produce a noticeable shift in the export time.

\subsection{Propagation of discrepancies along a chain of communities}
\label{sec:res_chain_random}

We examine successive inter-community crossings in a chain of ten communities.

We first consider a chain of cliques, in which internal infection spreads much faster than inter-community transmission. Even after nine consecutive crossings, the time distributions remained close (Figure~\ref{fig:chain_cliques}).

\begin{figure*}[!tbp]
\centering
% BEGIN EDITABLE FIGURE chain_cliques
\begin{tikzpicture}
\begin{groupplot}[paper axis,group style={group size=2 by 1,horizontal sep=1.8cm},width=.38\linewidth,height=4.6cm,scale only axis,xmin=0,xmax=6000,xtick={0,2000,4000,6000},xlabel={Time}]
\nextgroupplot[title={(a) Epidemic propagation},ylabel={Infected nodes},ymin=0,ymax=10300,ytick={0,2500,5000,7500,10000},scaled y ticks=false,legend pos=north west]
\addplot[draw=none,name path=fulllo,forget plot] table[x=time,y=lower]{figures_pgfplots/chain_full_trajectory.dat};
\addplot[draw=none,name path=fullhi,forget plot] table[x=time,y=upper]{figures_pgfplots/chain_full_trajectory.dat};
\addplot[fill=FullBlue,fill opacity=.16,draw=none,forget plot] fill between[of=fulllo and fullhi];
\addplot[FullBlue,thick] table[x=time,y=median]{figures_pgfplots/chain_full_trajectory.dat};
\addlegendentry{\Full{}}
\addplot[draw=none,name path=mmlo,forget plot] table[x=time,y=lower]{figures_pgfplots/chain_mm_trajectory.dat};
\addplot[draw=none,name path=mmhi,forget plot] table[x=time,y=upper]{figures_pgfplots/chain_mm_trajectory.dat};
\addplot[fill=ScaleOrange,fill opacity=.16,draw=none,forget plot] fill between[of=mmlo and mmhi];
\addplot[ScaleOrange,thick] table[x=time,y=median]{figures_pgfplots/chain_mm_trajectory.dat};
\addlegendentry{\TS{}}
\nextgroupplot[title={(b) Arrival in the last community},ylabel={Empirical CDF},ymin=0,ymax=1.02,ytick={0,.25,.5,.75,1}]
\addplot[FullBlue,thick,const plot] table[x=time,y=cdf]{figures_pgfplots/chain_full_cdf.dat};
\addplot[ScaleOrange,thick,const plot] table[x=time,y=cdf]{figures_pgfplots/chain_mm_cdf.dat};
\end{groupplot}
\end{tikzpicture}
% END EDITABLE FIGURE chain_cliques
\caption{%
Infection spread along a chain of ten cliques with 1000 nodes each ($k=1$, $\beta=0.003$, 400 runs of each model).
Left: median number of infected nodes over time and interquartile bands, which describe variability among runs rather than confidence intervals.
Right: empirical distribution functions of the arrival time at the last community, $T_9$. Four \Full{} values and two \TS{} values lie outside the displayed time range; all are included in the analysis.
}
\label{fig:chain_cliques}
\end{figure*}

The median arrival time at the last community was 2785 in \Full{} and 2877 in \TS{}. The normalised Wasserstein distance was 3.34\% and did not exceed the permutation threshold of 5.1\% ($p=0.32$). Thus, at this sample size, the test did not detect a statistically significant difference between the distributions.

For a clique with $n$ nodes and $\beta_{\mathrm{in}}=\beta_{\mathrm{out}}$, the ratio of the mean time to infect the entire clique to the mean transmission time through one active inter-community edge \michael{is $2H_{n-1}/n$, where $H_{n-1}=\sum_{m=1}^{n-1}1/m$ is a harmonic number; this ratio is only about $1.5\%$ for $n=1000$.}{is $2H_{n-1}/n$, which is only about $1.5\%$ for $n=1000$.}

Internal spread is therefore much faster than inter-community transmission, which explains the small discrepancies observed between the models. In addition, information about which particular nodes are bridges has relatively little effect on the crossing time.

A different pattern appears in chains of random communities with longer internal paths. \TS{} systematically predicts earlier infection arrival, but the size of this lead behaves differently in absolute and relative terms (Figure~\ref{fig:chain_heatmap}).

\begin{figure*}[!tbp]
\centering
% BEGIN EDITABLE FIGURE chain_heatmap
\begin{tikzpicture}
\begin{groupplot}[paper axis,group style={group size=2 by 2,horizontal sep=3.4cm,vertical sep=1.7cm},width=.31\linewidth,height=4.25cm,scale only axis]
\nextgroupplot[title={(a) $k=1$: median shift},xmin=-.5,xmax=8.5,ymin=-.5,ymax=5.5,xtick={0,1,2,3,4,5,6,7,8},xticklabels={1,2,3,4,5,6,7,8,9},ytick={0,1,2,3,4,5},yticklabels={1,2,3,4,5,6},xlabel={Community reached $c$},ylabel={Target internal ASPL},colormap name=paperdiv,point meta min=-14.28295105051232,point meta max=14.28295105051232,colorbar,colorbar style={width=3mm,ylabel={$b_c$ (\%)},ylabel style={font=\scriptsize},scaled y ticks=false},grid=none,tick align=outside]
\addplot[matrix plot*,mesh/cols=9,point meta=explicit] table[x=x,y=y,meta=value]{figures_pgfplots/chain_heatmap_k1_median.dat};
\pgfplotstableread{figures_pgfplots/chain_heatmap_k1_median.dat}\heatmapdata
\pgfplotstableforeachcolumnelement{value}\of\heatmapdata\as\cellvalue{
\pgfplotstablegetelem{\pgfplotstablerow}{x}\of\heatmapdata\edef\cellx{\pgfplotsretval}
\pgfplotstablegetelem{\pgfplotstablerow}{y}\of\heatmapdata\edef\celly{\pgfplotsretval}
\pgfplotstablegetelem{\pgfplotstablerow}{hatch}\of\heatmapdata
\ifnum\pgfplotsretval=1\relax\edef\drawcellhatch{\noexpand\path[pattern=north east lines,pattern color=gray!45] (axis cs:{\cellx-.5},{\celly-.5}) rectangle (axis cs:{\cellx+.5},{\celly+.5});}\drawcellhatch\fi
\pgfmathparse{abs(\cellvalue)>8.2841116092971454}
\ifdim\pgfmathresult pt>0pt\def\celltextcolor{white}\else\def\celltextcolor{black}\fi
\edef\drawcelllabel{\noexpand\node[font=\noexpand\ArrivalCellFont,text=\celltextcolor] at (axis cs:\cellx,\celly) {\noexpand\pgfmathprintnumber[assume math mode=true,fixed,precision=1,fixed zerofill,showpos=false]{\cellvalue}};}\drawcelllabel
}
\draw[white,line width=.4pt] (axis cs:0.5,-.5) -- (axis cs:0.5,5.5);
\draw[white,line width=.4pt] (axis cs:1.5,-.5) -- (axis cs:1.5,5.5);
\draw[white,line width=.4pt] (axis cs:2.5,-.5) -- (axis cs:2.5,5.5);
\draw[white,line width=.4pt] (axis cs:3.5,-.5) -- (axis cs:3.5,5.5);
\draw[white,line width=.4pt] (axis cs:4.5,-.5) -- (axis cs:4.5,5.5);
\draw[white,line width=.4pt] (axis cs:5.5,-.5) -- (axis cs:5.5,5.5);
\draw[white,line width=.4pt] (axis cs:6.5,-.5) -- (axis cs:6.5,5.5);
\draw[white,line width=.4pt] (axis cs:7.5,-.5) -- (axis cs:7.5,5.5);
\draw[white,line width=.4pt] (axis cs:-.5,0.5) -- (axis cs:8.5,0.5);
\draw[white,line width=.4pt] (axis cs:-.5,1.5) -- (axis cs:8.5,1.5);
\draw[white,line width=.4pt] (axis cs:-.5,2.5) -- (axis cs:8.5,2.5);
\draw[white,line width=.4pt] (axis cs:-.5,3.5) -- (axis cs:8.5,3.5);
\draw[white,line width=.4pt] (axis cs:-.5,4.5) -- (axis cs:8.5,4.5);
\nextgroupplot[title={(b) $k=1$: Wasserstein distance},xmin=-.5,xmax=8.5,ymin=-.5,ymax=5.5,xtick={0,1,2,3,4,5,6,7,8},xticklabels={1,2,3,4,5,6,7,8,9},ytick={0,1,2,3,4,5},yticklabels={1,2,3,4,5,6},xlabel={Community reached $c$},ylabel={Target internal ASPL},colormap name=paperseq,point meta min=0,point meta max=13.83377890219645,colorbar,colorbar style={width=3mm,ylabel={$e_W$ (\%)},ylabel style={font=\scriptsize},scaled y ticks=false},grid=none,tick align=outside]
\addplot[matrix plot*,mesh/cols=9,point meta=explicit] table[x=x,y=y,meta=value]{figures_pgfplots/chain_heatmap_k1_w1.dat};
\pgfplotstableread{figures_pgfplots/chain_heatmap_k1_w1.dat}\heatmapdata
\pgfplotstableforeachcolumnelement{value}\of\heatmapdata\as\cellvalue{
\pgfplotstablegetelem{\pgfplotstablerow}{x}\of\heatmapdata\edef\cellx{\pgfplotsretval}
\pgfplotstablegetelem{\pgfplotstablerow}{y}\of\heatmapdata\edef\celly{\pgfplotsretval}
\pgfplotstablegetelem{\pgfplotstablerow}{hatch}\of\heatmapdata
\ifnum\pgfplotsretval=1\relax\edef\drawcellhatch{\noexpand\path[pattern=north east lines,pattern color=gray!45] (axis cs:{\cellx-.5},{\celly-.5}) rectangle (axis cs:{\cellx+.5},{\celly+.5});}\drawcellhatch\fi
\pgfmathparse{abs(\cellvalue)>8.0235917632739397}
\ifdim\pgfmathresult pt>0pt\def\celltextcolor{white}\else\def\celltextcolor{black}\fi
\edef\drawcelllabel{\noexpand\node[font=\noexpand\ArrivalCellFont,text=\celltextcolor] at (axis cs:\cellx,\celly) {\noexpand\pgfmathprintnumber[assume math mode=true,fixed,precision=1,fixed zerofill,showpos=false]{\cellvalue}};}\drawcelllabel
}
\draw[white,line width=.4pt] (axis cs:0.5,-.5) -- (axis cs:0.5,5.5);
\draw[white,line width=.4pt] (axis cs:1.5,-.5) -- (axis cs:1.5,5.5);
\draw[white,line width=.4pt] (axis cs:2.5,-.5) -- (axis cs:2.5,5.5);
\draw[white,line width=.4pt] (axis cs:3.5,-.5) -- (axis cs:3.5,5.5);
\draw[white,line width=.4pt] (axis cs:4.5,-.5) -- (axis cs:4.5,5.5);
\draw[white,line width=.4pt] (axis cs:5.5,-.5) -- (axis cs:5.5,5.5);
\draw[white,line width=.4pt] (axis cs:6.5,-.5) -- (axis cs:6.5,5.5);
\draw[white,line width=.4pt] (axis cs:7.5,-.5) -- (axis cs:7.5,5.5);
\draw[white,line width=.4pt] (axis cs:-.5,0.5) -- (axis cs:8.5,0.5);
\draw[white,line width=.4pt] (axis cs:-.5,1.5) -- (axis cs:8.5,1.5);
\draw[white,line width=.4pt] (axis cs:-.5,2.5) -- (axis cs:8.5,2.5);
\draw[white,line width=.4pt] (axis cs:-.5,3.5) -- (axis cs:8.5,3.5);
\draw[white,line width=.4pt] (axis cs:-.5,4.5) -- (axis cs:8.5,4.5);
\nextgroupplot[title={(c) $k=5$: median shift},xmin=-.5,xmax=8.5,ymin=-.5,ymax=5.5,xtick={0,1,2,3,4,5,6,7,8},xticklabels={1,2,3,4,5,6,7,8,9},ytick={0,1,2,3,4,5},yticklabels={1,2,3,4,5,6},xlabel={Community reached $c$},ylabel={Target internal ASPL},colormap name=paperdiv,point meta min=-14.28295105051232,point meta max=14.28295105051232,colorbar=false,colorbar style={width=3mm,ylabel={$b_c$ (\%)},ylabel style={font=\scriptsize},scaled y ticks=false},grid=none,tick align=outside]
\addplot[matrix plot*,mesh/cols=9,point meta=explicit] table[x=x,y=y,meta=value]{figures_pgfplots/chain_heatmap_k5_median.dat};
\pgfplotstableread{figures_pgfplots/chain_heatmap_k5_median.dat}\heatmapdata
\pgfplotstableforeachcolumnelement{value}\of\heatmapdata\as\cellvalue{
\pgfplotstablegetelem{\pgfplotstablerow}{x}\of\heatmapdata\edef\cellx{\pgfplotsretval}
\pgfplotstablegetelem{\pgfplotstablerow}{y}\of\heatmapdata\edef\celly{\pgfplotsretval}
\pgfplotstablegetelem{\pgfplotstablerow}{hatch}\of\heatmapdata
\ifnum\pgfplotsretval=1\relax\edef\drawcellhatch{\noexpand\path[pattern=north east lines,pattern color=gray!45] (axis cs:{\cellx-.5},{\celly-.5}) rectangle (axis cs:{\cellx+.5},{\celly+.5});}\drawcellhatch\fi
\pgfmathparse{abs(\cellvalue)>8.2841116092971454}
\ifdim\pgfmathresult pt>0pt\def\celltextcolor{white}\else\def\celltextcolor{black}\fi
\edef\drawcelllabel{\noexpand\node[font=\noexpand\ArrivalCellFont,text=\celltextcolor] at (axis cs:\cellx,\celly) {\noexpand\pgfmathprintnumber[assume math mode=true,fixed,precision=1,fixed zerofill,showpos=false]{\cellvalue}};}\drawcelllabel
}
\draw[white,line width=.4pt] (axis cs:0.5,-.5) -- (axis cs:0.5,5.5);
\draw[white,line width=.4pt] (axis cs:1.5,-.5) -- (axis cs:1.5,5.5);
\draw[white,line width=.4pt] (axis cs:2.5,-.5) -- (axis cs:2.5,5.5);
\draw[white,line width=.4pt] (axis cs:3.5,-.5) -- (axis cs:3.5,5.5);
\draw[white,line width=.4pt] (axis cs:4.5,-.5) -- (axis cs:4.5,5.5);
\draw[white,line width=.4pt] (axis cs:5.5,-.5) -- (axis cs:5.5,5.5);
\draw[white,line width=.4pt] (axis cs:6.5,-.5) -- (axis cs:6.5,5.5);
\draw[white,line width=.4pt] (axis cs:7.5,-.5) -- (axis cs:7.5,5.5);
\draw[white,line width=.4pt] (axis cs:-.5,0.5) -- (axis cs:8.5,0.5);
\draw[white,line width=.4pt] (axis cs:-.5,1.5) -- (axis cs:8.5,1.5);
\draw[white,line width=.4pt] (axis cs:-.5,2.5) -- (axis cs:8.5,2.5);
\draw[white,line width=.4pt] (axis cs:-.5,3.5) -- (axis cs:8.5,3.5);
\draw[white,line width=.4pt] (axis cs:-.5,4.5) -- (axis cs:8.5,4.5);
\nextgroupplot[title={(d) $k=5$: Wasserstein distance},xmin=-.5,xmax=8.5,ymin=-.5,ymax=5.5,xtick={0,1,2,3,4,5,6,7,8},xticklabels={1,2,3,4,5,6,7,8,9},ytick={0,1,2,3,4,5},yticklabels={1,2,3,4,5,6},xlabel={Community reached $c$},ylabel={Target internal ASPL},colormap name=paperseq,point meta min=0,point meta max=13.83377890219645,colorbar=false,colorbar style={width=3mm,ylabel={$e_W$ (\%)},ylabel style={font=\scriptsize},scaled y ticks=false},grid=none,tick align=outside]
\addplot[matrix plot*,mesh/cols=9,point meta=explicit] table[x=x,y=y,meta=value]{figures_pgfplots/chain_heatmap_k5_w1.dat};
\pgfplotstableread{figures_pgfplots/chain_heatmap_k5_w1.dat}\heatmapdata
\pgfplotstableforeachcolumnelement{value}\of\heatmapdata\as\cellvalue{
\pgfplotstablegetelem{\pgfplotstablerow}{x}\of\heatmapdata\edef\cellx{\pgfplotsretval}
\pgfplotstablegetelem{\pgfplotstablerow}{y}\of\heatmapdata\edef\celly{\pgfplotsretval}
\pgfplotstablegetelem{\pgfplotstablerow}{hatch}\of\heatmapdata
\ifnum\pgfplotsretval=1\relax\edef\drawcellhatch{\noexpand\path[pattern=north east lines,pattern color=gray!45] (axis cs:{\cellx-.5},{\celly-.5}) rectangle (axis cs:{\cellx+.5},{\celly+.5});}\drawcellhatch\fi
\pgfmathparse{abs(\cellvalue)>8.0235917632739397}
\ifdim\pgfmathresult pt>0pt\def\celltextcolor{white}\else\def\celltextcolor{black}\fi
\edef\drawcelllabel{\noexpand\node[font=\noexpand\ArrivalCellFont,text=\celltextcolor] at (axis cs:\cellx,\celly) {\noexpand\pgfmathprintnumber[assume math mode=true,fixed,precision=1,fixed zerofill,showpos=false]{\cellvalue}};}\drawcelllabel
}
\draw[white,line width=.4pt] (axis cs:0.5,-.5) -- (axis cs:0.5,5.5);
\draw[white,line width=.4pt] (axis cs:1.5,-.5) -- (axis cs:1.5,5.5);
\draw[white,line width=.4pt] (axis cs:2.5,-.5) -- (axis cs:2.5,5.5);
\draw[white,line width=.4pt] (axis cs:3.5,-.5) -- (axis cs:3.5,5.5);
\draw[white,line width=.4pt] (axis cs:4.5,-.5) -- (axis cs:4.5,5.5);
\draw[white,line width=.4pt] (axis cs:5.5,-.5) -- (axis cs:5.5,5.5);
\draw[white,line width=.4pt] (axis cs:6.5,-.5) -- (axis cs:6.5,5.5);
\draw[white,line width=.4pt] (axis cs:7.5,-.5) -- (axis cs:7.5,5.5);
\draw[white,line width=.4pt] (axis cs:-.5,0.5) -- (axis cs:8.5,0.5);
\draw[white,line width=.4pt] (axis cs:-.5,1.5) -- (axis cs:8.5,1.5);
\draw[white,line width=.4pt] (axis cs:-.5,2.5) -- (axis cs:8.5,2.5);
\draw[white,line width=.4pt] (axis cs:-.5,3.5) -- (axis cs:8.5,3.5);
\draw[white,line width=.4pt] (axis cs:-.5,4.5) -- (axis cs:8.5,4.5);
\end{groupplot}
\end{tikzpicture}
% END EDITABLE FIGURE chain_heatmap
\caption{%
Change in discrepancies along a chain of ten communities with $n=1000$ nodes each.
Rows correspond to target ASPL values, and columns to the community reached ($c=1,\ldots,9$).
Results for $k=1$ are shown above and those for $k=5$ below.
Left: relative shift $b_c$ in the median arrival time; negative values mean that \TS{} predicts earlier arrival.
Right: normalised Wasserstein distance $\eW$.
Each configuration uses 400 paired runs. Uncertainty is not shown; results for terminal arrivals are presented in Figure~\ref{fig:accuracy_map}.
}
\label{fig:chain_heatmap}
\end{figure*}

At ASPL 6 and $k=1$, the median first-crossing time in \TS{} was 7.1\% lower, whereas the time to reach the last community was 13.8\% lower. For $k=5$, the corresponding values were 8.7\% and 10.1\%.

However, the relative discrepancy did not increase in proportion to the number of inter-community crossings. After the initial changes, it mostly stabilised or changed slowly. For example, at ASPL 5 and $k=1$, the median lead was 9.6\% for the first crossing and 8.7\% for arrival at the last community.

At the same time, the absolute difference between the median arrival times in this configuration increased from approximately 107 to 896 time units.

Successive inter-community crossings can accumulate a substantial absolute timing discrepancy even when its relative size remains approximately constant. This matters in forecasting problems, in which even a relatively small error can correspond to a substantial difference in the time at which infection reaches a distant community.

\subsection{Effects of community size and number of inter-community edges}
\label{sec:res_size}

At the same ASPL, the accuracy of \TS{} may depend on community size and the number of inter-community edges. We compare the arrival time at the last community, $T_9$, for different community sizes $n$, ASPL values, and numbers of inter-community edges $k$ (Figure~\ref{fig:accuracy_map}).

\begin{figure*}[!tbp]
\centering
% BEGIN EDITABLE FIGURE accuracy
\begin{tikzpicture}
\begin{groupplot}[paper axis,group style={group size=2 by 2,horizontal sep=3.4cm,vertical sep=1.7cm},width=.31\linewidth,height=4.0cm,scale only axis]
\nextgroupplot[title={(a) $k=1$: median shift},xmin=-.5,xmax=3.5,ymin=-.5,ymax=5.5,xtick={0,1,2,3},xticklabels={50,100,250,1000},ytick={0,1,2,3,4,5},yticklabels={1,2,3,4,5,6},xlabel={Community size $n$},ylabel={Target internal ASPL},colormap name=paperdiv,point meta min=-40.307303343994505,point meta max=40.307303343994505,colorbar,colorbar style={width=3mm,ylabel={$b_9$ (\%)},ylabel style={font=\scriptsize},scaled y ticks=false},grid=none,tick align=outside]
\addplot[matrix plot*,mesh/cols=4,point meta=explicit] table[x=x,y=y,meta=value]{figures_pgfplots/accuracy_k1_median.dat};
\pgfplotstableread{figures_pgfplots/accuracy_k1_median.dat}\heatmapdata
\pgfplotstableforeachcolumnelement{value}\of\heatmapdata\as\cellvalue{
\pgfplotstablegetelem{\pgfplotstablerow}{x}\of\heatmapdata\edef\cellx{\pgfplotsretval}
\pgfplotstablegetelem{\pgfplotstablerow}{y}\of\heatmapdata\edef\celly{\pgfplotsretval}
\pgfplotstablegetelem{\pgfplotstablerow}{hatch}\of\heatmapdata
\ifnum\pgfplotsretval=1\relax\edef\drawcellhatch{\noexpand\path[pattern=north east lines,pattern color=gray!45] (axis cs:{\cellx-.5},{\celly-.5}) rectangle (axis cs:{\cellx+.5},{\celly+.5});}\drawcellhatch\fi
\pgfmathparse{abs(\cellvalue)>23.378235939516813}
\ifdim\pgfmathresult pt>0pt\def\celltextcolor{white}\else\def\celltextcolor{black}\fi
\edef\drawcelllabel{\noexpand\node[font=\noexpand\scriptsize,text=\celltextcolor] at (axis cs:\cellx,\celly) {\noexpand\pgfmathprintnumber[fixed,precision=1,fixed zerofill,showpos=true]{\cellvalue}};}\drawcelllabel
}
\draw[white,line width=.4pt] (axis cs:0.5,-.5) -- (axis cs:0.5,5.5);
\draw[white,line width=.4pt] (axis cs:1.5,-.5) -- (axis cs:1.5,5.5);
\draw[white,line width=.4pt] (axis cs:2.5,-.5) -- (axis cs:2.5,5.5);
\draw[white,line width=.4pt] (axis cs:-.5,0.5) -- (axis cs:3.5,0.5);
\draw[white,line width=.4pt] (axis cs:-.5,1.5) -- (axis cs:3.5,1.5);
\draw[white,line width=.4pt] (axis cs:-.5,2.5) -- (axis cs:3.5,2.5);
\draw[white,line width=.4pt] (axis cs:-.5,3.5) -- (axis cs:3.5,3.5);
\draw[white,line width=.4pt] (axis cs:-.5,4.5) -- (axis cs:3.5,4.5);
\nextgroupplot[title={(b) $k=1$: Wasserstein distance},xmin=-.5,xmax=3.5,ymin=-.5,ymax=5.5,xtick={0,1,2,3},xticklabels={50,100,250,1000},ytick={0,1,2,3,4,5},yticklabels={1,2,3,4,5,6},xlabel={Community size $n$},ylabel={Target internal ASPL},colormap name=paperseq,point meta min=0,point meta max=40.88324313976721,colorbar,colorbar style={width=3mm,ylabel={$e_W$ (\%)},ylabel style={font=\scriptsize},scaled y ticks=false},grid=none,tick align=outside]
\addplot[matrix plot*,mesh/cols=4,point meta=explicit] table[x=x,y=y,meta=value]{figures_pgfplots/accuracy_k1_w1.dat};
\pgfplotstableread{figures_pgfplots/accuracy_k1_w1.dat}\heatmapdata
\pgfplotstableforeachcolumnelement{value}\of\heatmapdata\as\cellvalue{
\pgfplotstablegetelem{\pgfplotstablerow}{x}\of\heatmapdata\edef\cellx{\pgfplotsretval}
\pgfplotstablegetelem{\pgfplotstablerow}{y}\of\heatmapdata\edef\celly{\pgfplotsretval}
\pgfplotstablegetelem{\pgfplotstablerow}{hatch}\of\heatmapdata
\ifnum\pgfplotsretval=1\relax\edef\drawcellhatch{\noexpand\path[pattern=north east lines,pattern color=gray!45] (axis cs:{\cellx-.5},{\celly-.5}) rectangle (axis cs:{\cellx+.5},{\celly+.5});}\drawcellhatch\fi
\pgfmathparse{abs(\cellvalue)>23.712281021064982}
\ifdim\pgfmathresult pt>0pt\def\celltextcolor{white}\else\def\celltextcolor{black}\fi
\edef\drawcelllabel{\noexpand\node[font=\noexpand\scriptsize,text=\celltextcolor] at (axis cs:\cellx,\celly) {\noexpand\pgfmathprintnumber[fixed,precision=1,fixed zerofill,showpos=false]{\cellvalue}};}\drawcelllabel
}
\draw[white,line width=.4pt] (axis cs:0.5,-.5) -- (axis cs:0.5,5.5);
\draw[white,line width=.4pt] (axis cs:1.5,-.5) -- (axis cs:1.5,5.5);
\draw[white,line width=.4pt] (axis cs:2.5,-.5) -- (axis cs:2.5,5.5);
\draw[white,line width=.4pt] (axis cs:-.5,0.5) -- (axis cs:3.5,0.5);
\draw[white,line width=.4pt] (axis cs:-.5,1.5) -- (axis cs:3.5,1.5);
\draw[white,line width=.4pt] (axis cs:-.5,2.5) -- (axis cs:3.5,2.5);
\draw[white,line width=.4pt] (axis cs:-.5,3.5) -- (axis cs:3.5,3.5);
\draw[white,line width=.4pt] (axis cs:-.5,4.5) -- (axis cs:3.5,4.5);
\nextgroupplot[title={(c) $k=5$: median shift},xmin=-.5,xmax=3.5,ymin=-.5,ymax=5.5,xtick={0,1,2,3},xticklabels={50,100,250,1000},ytick={0,1,2,3,4,5},yticklabels={1,2,3,4,5,6},xlabel={Community size $n$},ylabel={Target internal ASPL},colormap name=paperdiv,point meta min=-40.307303343994505,point meta max=40.307303343994505,colorbar=false,colorbar style={width=3mm,ylabel={$b_9$ (\%)},ylabel style={font=\scriptsize},scaled y ticks=false},grid=none,tick align=outside]
\addplot[matrix plot*,mesh/cols=4,point meta=explicit] table[x=x,y=y,meta=value]{figures_pgfplots/accuracy_k5_median.dat};
\pgfplotstableread{figures_pgfplots/accuracy_k5_median.dat}\heatmapdata
\pgfplotstableforeachcolumnelement{value}\of\heatmapdata\as\cellvalue{
\pgfplotstablegetelem{\pgfplotstablerow}{x}\of\heatmapdata\edef\cellx{\pgfplotsretval}
\pgfplotstablegetelem{\pgfplotstablerow}{y}\of\heatmapdata\edef\celly{\pgfplotsretval}
\pgfplotstablegetelem{\pgfplotstablerow}{hatch}\of\heatmapdata
\ifnum\pgfplotsretval=1\relax\edef\drawcellhatch{\noexpand\path[pattern=north east lines,pattern color=gray!45] (axis cs:{\cellx-.5},{\celly-.5}) rectangle (axis cs:{\cellx+.5},{\celly+.5});}\drawcellhatch\fi
\pgfmathparse{abs(\cellvalue)>23.378235939516813}
\ifdim\pgfmathresult pt>0pt\def\celltextcolor{white}\else\def\celltextcolor{black}\fi
\edef\drawcelllabel{\noexpand\node[font=\noexpand\scriptsize,text=\celltextcolor] at (axis cs:\cellx,\celly) {\noexpand\pgfmathprintnumber[fixed,precision=1,fixed zerofill,showpos=true]{\cellvalue}};}\drawcelllabel
}
\draw[white,line width=.4pt] (axis cs:0.5,-.5) -- (axis cs:0.5,5.5);
\draw[white,line width=.4pt] (axis cs:1.5,-.5) -- (axis cs:1.5,5.5);
\draw[white,line width=.4pt] (axis cs:2.5,-.5) -- (axis cs:2.5,5.5);
\draw[white,line width=.4pt] (axis cs:-.5,0.5) -- (axis cs:3.5,0.5);
\draw[white,line width=.4pt] (axis cs:-.5,1.5) -- (axis cs:3.5,1.5);
\draw[white,line width=.4pt] (axis cs:-.5,2.5) -- (axis cs:3.5,2.5);
\draw[white,line width=.4pt] (axis cs:-.5,3.5) -- (axis cs:3.5,3.5);
\draw[white,line width=.4pt] (axis cs:-.5,4.5) -- (axis cs:3.5,4.5);
\nextgroupplot[title={(d) $k=5$: Wasserstein distance},xmin=-.5,xmax=3.5,ymin=-.5,ymax=5.5,xtick={0,1,2,3},xticklabels={50,100,250,1000},ytick={0,1,2,3,4,5},yticklabels={1,2,3,4,5,6},xlabel={Community size $n$},ylabel={Target internal ASPL},colormap name=paperseq,point meta min=0,point meta max=40.88324313976721,colorbar=false,colorbar style={width=3mm,ylabel={$e_W$ (\%)},ylabel style={font=\scriptsize},scaled y ticks=false},grid=none,tick align=outside]
\addplot[matrix plot*,mesh/cols=4,point meta=explicit] table[x=x,y=y,meta=value]{figures_pgfplots/accuracy_k5_w1.dat};
\pgfplotstableread{figures_pgfplots/accuracy_k5_w1.dat}\heatmapdata
\pgfplotstableforeachcolumnelement{value}\of\heatmapdata\as\cellvalue{
\pgfplotstablegetelem{\pgfplotstablerow}{x}\of\heatmapdata\edef\cellx{\pgfplotsretval}
\pgfplotstablegetelem{\pgfplotstablerow}{y}\of\heatmapdata\edef\celly{\pgfplotsretval}
\pgfplotstablegetelem{\pgfplotstablerow}{hatch}\of\heatmapdata
\ifnum\pgfplotsretval=1\relax\edef\drawcellhatch{\noexpand\path[pattern=north east lines,pattern color=gray!45] (axis cs:{\cellx-.5},{\celly-.5}) rectangle (axis cs:{\cellx+.5},{\celly+.5});}\drawcellhatch\fi
\pgfmathparse{abs(\cellvalue)>23.712281021064982}
\ifdim\pgfmathresult pt>0pt\def\celltextcolor{white}\else\def\celltextcolor{black}\fi
\edef\drawcelllabel{\noexpand\node[font=\noexpand\scriptsize,text=\celltextcolor] at (axis cs:\cellx,\celly) {\noexpand\pgfmathprintnumber[fixed,precision=1,fixed zerofill,showpos=false]{\cellvalue}};}\drawcelllabel
}
\draw[white,line width=.4pt] (axis cs:0.5,-.5) -- (axis cs:0.5,5.5);
\draw[white,line width=.4pt] (axis cs:1.5,-.5) -- (axis cs:1.5,5.5);
\draw[white,line width=.4pt] (axis cs:2.5,-.5) -- (axis cs:2.5,5.5);
\draw[white,line width=.4pt] (axis cs:-.5,0.5) -- (axis cs:3.5,0.5);
\draw[white,line width=.4pt] (axis cs:-.5,1.5) -- (axis cs:3.5,1.5);
\draw[white,line width=.4pt] (axis cs:-.5,2.5) -- (axis cs:3.5,2.5);
\draw[white,line width=.4pt] (axis cs:-.5,3.5) -- (axis cs:3.5,3.5);
\draw[white,line width=.4pt] (axis cs:-.5,4.5) -- (axis cs:3.5,4.5);
\end{groupplot}
\end{tikzpicture}
% END EDITABLE FIGURE accuracy
\caption{%
Discrepancies in the terminal arrival time $T_9$ for 48 configurations of ten-community chains.
The left panels show the relative median shift $b_9$ (\%), and the right panels show the normalised Wasserstein distance $\eW$ (\%). The top row corresponds to $k=1$ and the bottom row to $k=5$.
Within each panel, cell rows correspond to target ASPL values and columns to community size $n$.
The colour and number in a cell represent the same metric; the colour scales are shared between the two values of $k$ separately for each metric.
A negative median shift means earlier arrival in \TS{}.
Hatching in the right panels only marks configurations for which the permutation test for the Wasserstein distance did not detect a difference ($p\geq0.05$).
The tests were performed without correction for multiple comparisons.
Each configuration uses 400 paired runs.
}
\label{fig:accuracy_map}
\end{figure*}

In configurations with short internal paths ($\mathrm{ASPL}\leq2$), the discrepancies were mostly small. For 14 of these 16 configurations, the permutation test did not detect a statistically significant difference between the arrival-time distributions.

At $\mathrm{ASPL}\geq3$, the discrepancies became more pronounced. In 31 of the 32 configurations, the permutation test detected a difference between the distributions ($p<0.01$). These results were obtained without correction for multiple comparisons.

The magnitude of the discrepancy also depended strongly on community size. At fixed ASPL $\geq3$, the normalised Wasserstein distance decreased monotonically with increasing $n$ for both values of $k$.

For example, at ASPL 6 and $k=1$, it was 40.9\% for $n=50$ but only 13.6\% for $n=1000$. For $k=5$, the corresponding values were 32.6\% and 10.8\%.

Increasing the number of inter-community edges from one to five also generally reduced the discrepancies in sparse communities, although it did not remove them completely.

The fraction of \TS{} inter-community crossings that occurred before infection of the first real bridge in the source also depended strongly on structure: it ranged from 0.1--2.3\% at ASPL 1 to 31--51\% at ASPL 6.

Discrepancies between the models generally increase as internal paths become longer. However, at the same ASPL, their magnitude can differ substantially with community size and the number of inter-community edges. ASPL alone is insufficient for assessing accuracy.

\subsection{Relationship between discrepancies and the ratio of time scales}
\label{sec:res_descriptor}

We assess the relationship between discrepancies and temporal heterogeneity of internal infection using the measure $\rho=\beta_{\mathrm{out}}\sigma_T$ defined in Section~\ref{sec:descriptor}.

Figure~\ref{fig:descriptor}(a,b) shows the normalised Wasserstein distance $\eW$ and median shift $b_9$ as functions of $\rho$ for all 48 chain configurations.

\begin{figure*}[!tbp]
\centering
% BEGIN EDITABLE FIGURE descriptor
\begin{tikzpicture}
\begin{groupplot}[paper axis,group style={group size=3 by 1,horizontal sep=1.7cm},width=.24\linewidth,height=4.1cm,scale only axis]
\nextgroupplot[xmode=log,xlabel={$\rho=\beta\sigma_T$},ylabel={$e_W$ of $T_9$ (\%)},title={(a) Distributional discrepancy}]
\addplot[Size50,mark=*,mark size=1.7pt,mark options={fill=Size50}] table[x=plotrho,y=rel_w1_pct]{figures_pgfplots/descriptor_n50_k1.dat};
\addplot[gray!65,only marks,mark=-,mark size=2pt] table[x=plotrho,y=rel_w1_null95_pct]{figures_pgfplots/descriptor_n50_k1.dat};
\addplot[Size50,mark=triangle,mark size=1.7pt,mark options={fill=white}] table[x=plotrho,y=rel_w1_pct]{figures_pgfplots/descriptor_n50_k5.dat};
\addplot[gray!65,only marks,mark=-,mark size=2pt] table[x=plotrho,y=rel_w1_null95_pct]{figures_pgfplots/descriptor_n50_k5.dat};
\addplot[Size100,mark=*,mark size=1.7pt,mark options={fill=Size100}] table[x=plotrho,y=rel_w1_pct]{figures_pgfplots/descriptor_n100_k1.dat};
\addplot[gray!65,only marks,mark=-,mark size=2pt] table[x=plotrho,y=rel_w1_null95_pct]{figures_pgfplots/descriptor_n100_k1.dat};
\addplot[Size100,mark=triangle,mark size=1.7pt,mark options={fill=white}] table[x=plotrho,y=rel_w1_pct]{figures_pgfplots/descriptor_n100_k5.dat};
\addplot[gray!65,only marks,mark=-,mark size=2pt] table[x=plotrho,y=rel_w1_null95_pct]{figures_pgfplots/descriptor_n100_k5.dat};
\addplot[Size250,mark=*,mark size=1.7pt,mark options={fill=Size250}] table[x=plotrho,y=rel_w1_pct]{figures_pgfplots/descriptor_n250_k1.dat};
\addplot[gray!65,only marks,mark=-,mark size=2pt] table[x=plotrho,y=rel_w1_null95_pct]{figures_pgfplots/descriptor_n250_k1.dat};
\addplot[Size250,mark=triangle,mark size=1.7pt,mark options={fill=white}] table[x=plotrho,y=rel_w1_pct]{figures_pgfplots/descriptor_n250_k5.dat};
\addplot[gray!65,only marks,mark=-,mark size=2pt] table[x=plotrho,y=rel_w1_null95_pct]{figures_pgfplots/descriptor_n250_k5.dat};
\addplot[Size1000,mark=*,mark size=1.7pt,mark options={fill=Size1000}] table[x=plotrho,y=rel_w1_pct]{figures_pgfplots/descriptor_n1000_k1.dat};
\addplot[gray!65,only marks,mark=-,mark size=2pt] table[x=plotrho,y=rel_w1_null95_pct]{figures_pgfplots/descriptor_n1000_k1.dat};
\addplot[Size1000,mark=triangle,mark size=1.7pt,mark options={fill=white}] table[x=plotrho,y=rel_w1_pct]{figures_pgfplots/descriptor_n1000_k5.dat};
\addplot[gray!65,only marks,mark=-,mark size=2pt] table[x=plotrho,y=rel_w1_null95_pct]{figures_pgfplots/descriptor_n1000_k5.dat};
\nextgroupplot[xmode=log,xlabel={$\rho=\beta\sigma_T$},ylabel={Median shift $b_9$ (\%)},title={(b) Signed median shift}]
\addplot[gray,no marks] coordinates {(0.0018131616755314697,0) (2.7344524800897014,0)};
\addplot[Size50,mark=*,mark size=1.7pt,mark options={fill=Size50}] table[x=plotrho,y=median_shift_pct]{figures_pgfplots/descriptor_n50_k1.dat};
\addplot[Size50,mark=triangle,mark size=1.7pt,mark options={fill=white}] table[x=plotrho,y=median_shift_pct]{figures_pgfplots/descriptor_n50_k5.dat};
\addplot[Size100,mark=*,mark size=1.7pt,mark options={fill=Size100}] table[x=plotrho,y=median_shift_pct]{figures_pgfplots/descriptor_n100_k1.dat};
\addplot[Size100,mark=triangle,mark size=1.7pt,mark options={fill=white}] table[x=plotrho,y=median_shift_pct]{figures_pgfplots/descriptor_n100_k5.dat};
\addplot[Size250,mark=*,mark size=1.7pt,mark options={fill=Size250}] table[x=plotrho,y=median_shift_pct]{figures_pgfplots/descriptor_n250_k1.dat};
\addplot[Size250,mark=triangle,mark size=1.7pt,mark options={fill=white}] table[x=plotrho,y=median_shift_pct]{figures_pgfplots/descriptor_n250_k5.dat};
\addplot[Size1000,mark=*,mark size=1.7pt,mark options={fill=Size1000}] table[x=plotrho,y=median_shift_pct]{figures_pgfplots/descriptor_n1000_k1.dat};
\addplot[Size1000,mark=triangle,mark size=1.7pt,mark options={fill=white}] table[x=plotrho,y=median_shift_pct]{figures_pgfplots/descriptor_n1000_k5.dat};
\nextgroupplot[title={(c) First vs later crossings},xmin=-47,xmax=5,ymin=-47,ymax=5,xtick={-40,-20,0},ytick={-40,-20,0},xlabel={First crossing: mean shift (\%)},ylabel={Crossings 2--9: mean shift (\%)}]
\addplot[gray,no marks] coordinates {(-47,-47) (5,5)};
\addplot[only marks,Size50,mark=*,mark size=1.7pt,mark options={fill=Size50},error bars/.cd,x dir=both,x explicit,error bar style={line width=.3pt},error mark options={mark size=0pt}] table[x=first_cross_shift_pct,y=later_cross_shift_pct,x error minus=errlo,x error plus=errhi]{figures_pgfplots/descriptor_cross_n50_k1.dat};
\addplot[only marks,Size50,mark=triangle,mark size=1.7pt,mark options={fill=white},error bars/.cd,x dir=both,x explicit,error bar style={line width=.3pt},error mark options={mark size=0pt}] table[x=first_cross_shift_pct,y=later_cross_shift_pct,x error minus=errlo,x error plus=errhi]{figures_pgfplots/descriptor_cross_n50_k5.dat};
\addplot[only marks,Size100,mark=*,mark size=1.7pt,mark options={fill=Size100},error bars/.cd,x dir=both,x explicit,error bar style={line width=.3pt},error mark options={mark size=0pt}] table[x=first_cross_shift_pct,y=later_cross_shift_pct,x error minus=errlo,x error plus=errhi]{figures_pgfplots/descriptor_cross_n100_k1.dat};
\addplot[only marks,Size100,mark=triangle,mark size=1.7pt,mark options={fill=white},error bars/.cd,x dir=both,x explicit,error bar style={line width=.3pt},error mark options={mark size=0pt}] table[x=first_cross_shift_pct,y=later_cross_shift_pct,x error minus=errlo,x error plus=errhi]{figures_pgfplots/descriptor_cross_n100_k5.dat};
\addplot[only marks,Size250,mark=*,mark size=1.7pt,mark options={fill=Size250},error bars/.cd,x dir=both,x explicit,error bar style={line width=.3pt},error mark options={mark size=0pt}] table[x=first_cross_shift_pct,y=later_cross_shift_pct,x error minus=errlo,x error plus=errhi]{figures_pgfplots/descriptor_cross_n250_k1.dat};
\addplot[only marks,Size250,mark=triangle,mark size=1.7pt,mark options={fill=white},error bars/.cd,x dir=both,x explicit,error bar style={line width=.3pt},error mark options={mark size=0pt}] table[x=first_cross_shift_pct,y=later_cross_shift_pct,x error minus=errlo,x error plus=errhi]{figures_pgfplots/descriptor_cross_n250_k5.dat};
\addplot[only marks,Size1000,mark=*,mark size=1.7pt,mark options={fill=Size1000},error bars/.cd,x dir=both,x explicit,error bar style={line width=.3pt},error mark options={mark size=0pt}] table[x=first_cross_shift_pct,y=later_cross_shift_pct,x error minus=errlo,x error plus=errhi]{figures_pgfplots/descriptor_cross_n1000_k1.dat};
\addplot[only marks,Size1000,mark=triangle,mark size=1.7pt,mark options={fill=white},error bars/.cd,x dir=both,x explicit,error bar style={line width=.3pt},error mark options={mark size=0pt}] table[x=first_cross_shift_pct,y=later_cross_shift_pct,x error minus=errlo,x error plus=errhi]{figures_pgfplots/descriptor_cross_n1000_k5.dat};
\node[anchor=north west,font=\tiny,align=left,text=gray] at (axis cs:-45,3) {ASPL $\geq3$\\Below diagonal: later\\crossings more biased};
\end{groupplot}
\node[anchor=north,font=\scriptsize] at ([yshift=-1.25cm]group c2r1.south) {\textcolor{Size50}{$\blacksquare$} $n=50$\quad\textcolor{Size100}{$\blacksquare$} $n=100$\quad\textcolor{Size250}{$\blacksquare$} $n=250$\quad\textcolor{Size1000}{$\blacksquare$} $n=1000$\quad$\bullet$ $k=1$\quad$\triangle$ $k=5$\quad\textcolor{gray}{---} null 95\%};
\end{tikzpicture}
% END EDITABLE FIGURE descriptor
\caption{%
Relationship between discrepancies in \Full{} and \TS{} and the time-scale ratio $\rho=\beta_{\mathrm{out}}\sigma_T$ for 48 chain configurations.
(a) Normalised Wasserstein distance $\eW$ for the terminal arrival time $T_9$.
(b) Relative shift in its median, $b_9$.
Colours denote community size; filled circles correspond to $k=1$, and open triangles to $k=5$.
Lines connect configurations with ASPL from 1 to 6 at fixed $n$ and $k$. Grey marks in (a) show the 95th percentile of the permutation distribution of the distance for the corresponding configuration.
(c) Comparison of relative shifts in the mean duration of the first crossing and crossings 2--9 for $\mathrm{ASPL}\geq3$.
Error bars show 95\% block-bootstrap intervals for the shift in the first crossing.
}
\label{fig:descriptor}
\end{figure*}

The measure $\rho$ had a stronger rank association with $\eW$ than target ASPL. The Spearman correlation between $\rho$ and $\eW$ was 0.91 across all configurations and 0.98 for configurations with $\mathrm{ASPL}\geq3$. The corresponding values for ASPL were 0.82 and 0.66.

For $\rho<0.25$, the distance $\eW$ \michael{did not exceed 4.53\%}{did not exceed 4.5\%} in any of the 18 configurations. For 15 of them, the permutation test did not detect a difference between the distributions. By contrast, for $\rho\geq0.25$, the observed distance exceeded the permutation threshold in all 30 configurations.

For $\rho\geq0.3$, the results were approximately described by the empirical relationship
\begin{equation}
\michael{\eW \approx 15\rho\qquad(\eW\text{ expressed in percent}).}{\eW \approx 15\rho\,\%.}
\end{equation}
In this range, we also observed $b_9\approx-\eW$: the distance between the distributions mainly corresponded to a shift towards earlier arrival times in \TS{}.

\michael{In these benchmarks, a small spread of internal infection times relative to the single-edge transmission time was associated with small distributional discrepancies. The observed thresholds and fitted relationship are descriptive of these configurations, not universal error guarantees.}{When the spread of internal infection times is small relative to the expected inter-community transmission time, aggregation of bridge nodes has a comparatively small effect on the time distributions.}

\section{Discussion}
\label{sec:discussion}

The main advantage of the proposed two-scale approach is that it combines detailed stochastic simulation within communities with an aggregate representation of the connections between them. This is suitable for problems in which contact-network data are available at different levels of detail. For example, the internal structure of communities can be reconstructed or generated from known characteristics, whereas inter-community connections may be represented only by aggregate movement flows obtained from transport or migration data. Unlike full-network simulation, \TS{} does not require the nodes involved in inter-community contacts to be specified explicitly. At the same time, unlike fully aggregated models, it preserves individual infection dynamics and the effect of internal topology.

The experiments showed that this simplification can preserve high accuracy in the temporal characteristics of infection spread. In communities with fast internal spread, the arrival-time distributions from \Full{} and \TS{} were close, and this agreement persisted even after nine consecutive inter-community crossings. Thus, the absence of information about specific bridge nodes does not necessarily cause a substantial error in the overall propagation dynamics. This makes aggregate interfaces a promising way to model networks of communities when the main quantities of interest are the times at which communities are reached and their distributions.

The accuracy of this approximation can be explained by the ratio of time scales. If infection spreads rapidly among the nodes of a community relative to the waiting time for inter-community transmission, the specific locations of the bridges have a comparatively small effect on the export time. With slower and more heterogeneous internal spread, however, the states of individual bridge nodes may remain important for a substantial part of the waiting time for a crossing. This is why structural characteristics, including ASPL, do not fully explain the observed discrepancies. The stronger relationship between accuracy and the measure $\rho=\beta_{\mathrm{out}}\sigma_T$ supports an explanation based on temporal heterogeneity in internal infection. This suggests that the suitability of aggregation can be assessed from both network topology and process dynamics.

A possible direction for further research is to increase the model's level of detail by incorporating local community characteristics, such as ASPL, eigenvalues, or other descriptors, and to evaluate whether these additions reduce the discrepancy between the methods.

For simulations of spread through many communities, the way in which discrepancies accumulate is also important. The experiments showed that the absolute difference between arrival times can increase along the chain, whereas the relative shift often stabilises or changes slowly. Thus, the error of the aggregate simulation does not necessarily grow in proportion to the number of inter-community crossings. At the same time, the stronger bias in the total duration of later crossings than in the first crossing indicates possible additional effects after infection is imported. One explanation is the occurrence of surplus imports, which can create additional foci of infection and accelerate subsequent spread. The separate contribution of this mechanism requires further investigation.

The theoretical analysis clearly identifies the approximation implemented by the two-scale approach. The ideal process is equivalent to exact simulation on a surrogate network with uniformly distributed inter-community connections. This provides a mathematical basis for the aggregate representation and explains why its results can differ from those of full-network simulation. In particular, similarity between time distributions does not guarantee reproduction of specific infection routes. The appropriate level of detail should therefore be determined by the aim of the study: an aggregate representation may be sufficient for estimating temporal characteristics, whereas analysis of physical transmission paths requires information about the corresponding contacts.

Further possibilities arise from the computational structure of the method. Independent local simulators can run in parallel between synchronisation points, and individual communities can have different internal structures and simulation parameters. In addition, replacing explicit inter-community edges with aggregate weights can reduce the amount of contact data. \michael{However, the actual gains in speed and memory depend on the network structure and the cost of synchronisation. The experiments use static synthetic SI networks and one synchronisation step ($\tau=1$); without a step-size convergence study, the measured discrepancy cannot be assigned entirely to interface aggregation.

\textbf{Outlook.} A natural next step is evaluation on openly available empirical contact networks, for example the SocioPatterns high-school contact data \cite{mastrandrea2015} and primary-school contact data \cite{stehle2011,gemmetto2014}. Their class labels provide community partitions, and measured contacts permit comparison of a full-network reference with an interface retaining only $W_{ij}$. A first study could use explicitly defined static contact aggregates; retaining the recorded contact times would require a temporal-network extension. Proper evaluation must address observation windows, contact-weight calibration, incomplete participation, and sensitivity to disease parameters and community definitions. These tasks warrant a dedicated follow-up paper beyond the present methodological study. Such data test fidelity on observed contact structures; validation against observed epidemic outcomes would require additional evidence.}{However, the actual gains in speed and memory depend on the network structure and the cost of synchronisation.}

Geographically explicit synthetic networks provide another setting for evaluation \cite{jiang2024synthetic}. Residential locations can define city or regional communities, with individual connections retained for the full-network reference and aggregated into $W_{ij}$ for the two-scale model. Where individual connections between communities are unavailable, local contact networks can be constructed from demographic data and linked using weights informed by commuting flows \cite{tulchinsky2024networks,tizzoni2014mobility}. This requires an explicit assumption about how mobility translates into contact intensity.

\section{Conclusions}
\label{sec:conclusions}

This study investigated a two-scale method for stochastic simulation of infection spread in networks of communities. The method combines detailed Gillespie simulation of internal contact networks with aggregate inter-community intensities. It therefore allows local contact networks to be used together with aggregate data on connections between communities, without requiring every inter-community contact to be represented explicitly. Its modular structure also allows local simulators to be run in parallel.

We proved that the ideal two-scale process is equivalent to exact simulation on a surrogate network with uniform inter-community connections. We analysed the mechanisms that produce discrepancies between aggregate and full-network simulation, including the consequences of losing information about specific bridge nodes. \michael{We also proved that, in the single-source--leaf setting of Proposition~\ref{prop:ordering}, uniform random bridge placement yields a first infection export that is stochastically no later in the ideal aggregate model.}{We also proved that, when bridge nodes are placed uniformly at random, aggregation yields a first infection export that is stochastically no later in the ideal model.}

The experiments confirmed that \TS{} can reproduce infection-spread time distributions close to those of \Full{} when infection spreads rapidly within communities. This agreement persisted even after nine consecutive inter-community crossings. The absolute discrepancy between arrival times can accumulate along a chain, whereas the relative discrepancy often stabilises or changes slowly.

We examined how internal structure, community size, and the number of inter-community edges affect the accuracy of aggregate simulation. The time-scale ratio $\rho=\beta_{\mathrm{out}}\sigma_T$ had a stronger observed relationship with the discrepancies than ASPL and is a promising characteristic for identifying conditions under which the method is applicable.

In our tests, the smallest discrepancies occurred when infection spread rapidly within communities. In all 18 configurations with $\rho<0.25$, the normalised Wasserstein distance \michael{did not exceed 4.53\%}{did not exceed 4.5\%}. At the same time, close time distributions did not guarantee preservation of admissible transmission routes.

% The appendices restore material already present in the complete TwoScale draft.
% Recovered from main-SMPT-twoscale.tex; scientific content retained.
\appendices

\section{Clique formulas and verification}
\label{app:clique}

Consider an $n$-node clique with SI dynamics, per-edge rate $\beta$, a uniformly placed seed and a uniformly placed set of $k$ bridge nodes. When $m$ nodes are infected, the next internal infection occurs at rate $\beta m(n-m)$, independently of which nodes are infected, and the infected set is a uniform $m$-subset. The probability that none of the first $m$ infected nodes is a bridge is $\binom{n-k}{m}/\binom{n}{m}$, hence
\begin{equation}
\label{eq:clique_TB}
\mathbb{E}[T_B]=\sum_{m=1}^{n-k}\frac{\binom{n-k}{m}}{\binom{n}{m}}\frac{1}{\beta m(n-m)},
\end{equation}
For one bridge this simplifies to
\begin{equation}
\mathbb{E}[T_B]=\frac{H_{n-1}}{\beta n}\ \ (k=1),
\end{equation}
where $H_{n-1}$ is the harmonic number; the case of a bridge seed ($T_B=0$) is included. For $k=1$, the export time in \Full{} is $T_B$ plus an independent $\Exp(\beta)$ waiting time. The expected time to infect the whole clique is $2H_{n-1}/(\beta n)$. Table~\ref{tab:verification} compares~\eqref{eq:clique_TB} with the first community of the clique chains.

\begin{table}[!htbp]
\centering
\caption{Verification of the reference simulator on the first community of the clique chains ($\beta=0.003$, 400 runs per cell): mean bridge-activation time and, for $k=1$, mean wait after activation (exact value $1/\beta=333.3$). Values are mean $\pm$ standard error.}
\label{tab:verification}
\small
\setlength{\tabcolsep}{3pt}
\begin{tabular}{rrrrr}
\toprule
$n$ & $k$ & Exact $\mathbb{E}[T_B]$ & \Full{} $T_B$ & \Full{} wait\\
\midrule
50 & 1 & 29.86 & $29.72\pm0.71$ & $354.0\pm16.9$\\
100 & 1 & 17.26 & $17.05\pm0.38$ & $326.5\pm17.0$\\
250 & 1 & 8.13 & $7.84\pm0.15$ & $347.4\pm16.6$\\
1000 & 1 & 2.49 & $2.51\pm0.04$ & $353.8\pm17.9$\\
50 & 5 & 15.97 & $15.96\pm0.55$ & --\\
100 & 5 & 10.31 & $9.64\pm0.29$ & --\\
250 & 5 & 5.35 & $5.16\pm0.12$ & --\\
1000 & 5 & 1.80 & $1.77\pm0.03$ & --\\
\bottomrule
\end{tabular}
\end{table}

\yulian{}{\section{Sampling floor of the Wasserstein distance}
\label{app:floor}

Let $F_R$ and $G_R$ be empirical CDFs of two independent samples of size $R$ from the same continuous law $F$. Then $W_1=\int|F_R-G_R|\,dx$, and $\sqrt{R}\,(F_R-G_R)$ is approximately $\sqrt{2}$ times a Brownian bridge evaluated at $F(x)$ \cite{delbarrio1999}. Since $\mathbb{E}|B(u)|=\sqrt{2u(1-u)/\pi}$ for a Brownian bridge $B$,
\[
\mathbb{E}[W_1]\approx\frac{2}{\sqrt{\pi R}}\int\sqrt{F(x)\bigl(1-F(x)\bigr)}\,dx .
\]
For an exponential law with mean $\mu$, the integral equals $\mu\pi/2$, so that $\mathbb{E}[W_1]/\mu\approx\sqrt{\pi/R}$, which is 8.9\% for $R=400$ and 1.8\% for $R=10\,000$. Arrival times that are sums of several crossings are less dispersed relative to their mean, and their floor is correspondingly lower; the permutation floors of the chain configurations ranged from 1.6\% to 5.9\%.

\section{Reproducibility details}
\label{app:repro}

Seeds are derived deterministically from the experiment parameters. Each \TS{} engine draws from a private pseudorandom stream; midpoint probes copy the stream state of their engine and do not advance it. Two \Full{} runs ($n=50$, ASPL 6, $k=1$) had not reached the tenth community by $T=30\,000$. They were rerun to $T=40\,000$ with the same graphs, initial nodes and event seeds. Their event histories up to 30\,000 were verified to match the saved runs, and their terminal arrivals occurred at about 31\,059 and 31\,007. No horizon value was substituted for an unobserved event. The permutation and bootstrap analyses of the chain configurations, the descriptor $\sigma_T$ and Figures~\ref{fig:accuracy_map}--\ref{fig:descriptor} were computed from the saved runs by the scripts \path{extrascripts/analysis/smpt_twoscale_paper_metrics.py} and \path{extrascripts/visualization/smpt_twoscale_paper_figures.py}.}

\section*{Code availability}
The source code of the simulators is available at
\url{https://github.com/YulianKuryliak/Two-Layer-CTMC}.

% Author-specific declarations below require completion before submission.
\section*{Data availability}
\michael{The simulation data are available at
\url{https://doi.org/10.5281/zenodo.23591882}.}{The data is available at
\url{10.5281/zenodo.23591882}.}

\section*{Funding}
This research did not receive any specific funding.

\section*{CRediT authorship contribution statement}

\textbf{Yulian Kuryliak:} Conceptualization, Methodology, Software,
Investigation, Formal analysis, Visualization, Writing -- original draft.
\textbf{Michael T. M. Emmerich:} Conceptualization, Methodology, Supervision,
Writing -- review \& editing.

\section*{Declaration of competing interest}
The authors declare that they have no known competing financial interests or
personal relationships that could have appeared to influence the work reported
in this article.

\section*{Acknowledgment}
\michael{The scientific direction and ideas originated solely with the authors. Anthropic Claude and OpenAI Codex were used for programming, plotting assistance, language editing and reformulation. The authors take full responsibility for all results, references, and     statements in this article.}{During the preparation of this manuscript, the first author used OpenAI ChatGPT and Codex, and Anthropic Claude to assist with drafting, language editing, and structuring the manuscript across Sections I–VIII and Appendices A–C. These tools were also used interactively to discuss and refine mathematical arguments, support experimental design and statistical analysis, and assist with the development of simulation and analysis code.

The authors directed the research, made the final scientific and methodological decisions, and reviewed and approved the resulting manuscript.}

\bibliographystyle{IEEEtran}
% Balance the final columns with the native IEEE reference trigger.
\IEEEtriggeratref{5}
\michael{\bibliography{references-IEEE-michael}}{\bibliography{references_fixed,references_missing}}
% Keep the final page compact; the class otherwise stretches the author gap.
\patchcmd{\IEEEbiographynophoto}
  {\vskip 4\baselineskip plus 1fil minus 0\baselineskip}
  {\vskip 2\baselineskip}
  {}{\PackageError{manuscript}{Biography spacing patch failed}{Check the IEEE Access class.}}
\raggedbottom
\begin{IEEEbiographynophoto}{Yulian Kuryliak}
is currently pursuing a Ph.D. in System Analysis at Lviv Polytechnic National University, Ukraine, where he also workes as a Teaching Assistant.
His research interests include stochastic modeling, complex networks, multiscale simulation, and the modeling of dynamical processes on networks, particularly epidemic spreading.
\end{IEEEbiographynophoto}

\begin{IEEEbiographynophoto}{Michael T.M. Emmerich}
Michael Emmerich received the Doctor of Natural Sciences degree (magna cum laude) from the Technical University of Dortmund, in 2005.
With more than 20 years of experience in academia—having worked at FOM, Leiden University, The Netherlands, IST Lisbon, Portugal, and RWTH Aachen, Germany, and in industry, including as a former Lead AI Scientist at silo.ai, he has established himself as a leading figure in his field. Currently, he is a Professor in Intelligent Computing with the University of Jyväskylä, Finland. He is a key contributor to the development of widely recognized algorithms in multi-objective optimization, such as SMS-EMOA, Hypervolume Indicator Gradient Ascent, Optimally Weighted Cluster Kriging, and Hypervolume Expected Improvement, and worked on different application domains, such as drug discovery, stochastic simulation of biological processes, and complex networks.
\end{IEEEbiographynophoto}

\EOD
\end{document}
